\documentclass[11pt]{article}

\usepackage[T1]{fontenc}
\usepackage[utf8]{inputenc}
\usepackage[english]{babel}
\usepackage[margin=1in]{geometry}
\usepackage{amsmath,amssymb,amsthm}
\usepackage{booktabs}
\usepackage{enumitem}
\usepackage{xcolor}
\usepackage{hyperref}
\usepackage{listings}
\usepackage{cleveref}

\emergencystretch=3em

\hypersetup{
  colorlinks=true,
  linkcolor=blue!60!black,
  citecolor=blue!60!black,
  urlcolor=blue!60!black
}

\theoremstyle{plain}
\newtheorem{theorem}{Theorem}[section]
\newtheorem{lemma}[theorem]{Lemma}
\newtheorem{proposition}[theorem]{Proposition}
\newtheorem{corollary}[theorem]{Corollary}
\newtheorem{conjecture}[theorem]{Conjecture}

\theoremstyle{definition}
\newtheorem{definition}[theorem]{Definition}
\newtheorem{example}[theorem]{Example}
\newtheorem{remark}[theorem]{Remark}

\newcommand{\N}{\mathbb{N}}
\newcommand{\Z}{\mathbb{Z}}
\newcommand{\Q}{\mathbb{Q}}
\newcommand{\R}{\mathbb{R}}
\newcommand{\vtwo}[1]{v_2\!\left(#1\right)}
\newcommand{\oddpart}{\mathrm{oddpart}}
\newcommand{\Tstep}{T}
\newcommand{\Orbit}{\mathrm{Orbit}}
\newcommand{\OrbitCycle}{\mathrm{OrbitCycle}}
\newcommand{\wordOrbit}{\mathrm{wordOrbit}}
\newcommand{\wordA}{A}
\newcommand{\wordWeight}{W}
\newcommand{\Pmi}{\mathrm{PMI}}
\newcommand{\FullOrbit}{\mathrm{FullOrbitFrom7}}
\newcommand{\GeneralOrbit}{\mathrm{GeneralOrbitFrom7}}
\newcommand{\ResetWindow}{\mathrm{ResetWindowReachability}}

\newcommand{\statusmath}{\textbf{Math: proven}}
\newcommand{\statuslean}{\textbf{Lean: compiled}}
\newcommand{\statuspending}{\textbf{Lean: pending}}

\title{The $5x+1$ orbit of $7$ is unbounded:\\ a string-flow proof with Lean verification}
\author{StringFlow}
\date{\today}

\begin{document}

\sloppy

\maketitle

\begin{abstract}
We study the accelerated $5x+1$ map
$T(x)=(5x+1)/2^{\vtwo{5x+1}}$ and prove that the orbit of $7$ is
unbounded.  The proof is organised around an exact string-flow
accounting of the orbit: every step is represented by a word of
2-adic weights, every block satisfies an exact prefix identity, and the
remaining obstruction is reduced to a corrected decisive window bound
that carries genuine orbit reachability.  The mathematical layer of the
unified core is closed; the corresponding Lean statement
\texttt{unified\_core\_final\_no\_hge} is the single remaining formalisation
gap and is explicitly marked as pending in this paper.
\end{abstract}

\tableofcontents

\section{Introduction}\label{sec:intro}

The $5x+1$ problem asks for the behaviour of the map
\begin{equation}\label{eq:step}
  T(x)=\frac{5x+1}{2^{\vtwo{5x+1}}},
\end{equation}
where $\vtwo{n}$ denotes the exponent of $2$ in the factorisation of a
positive integer $n$.  The map is the $5x+1$ analogue of the
Collatz map: odd multipliers increase the size of the state, while the
removal of powers of two may decrease it.  For background on exact
cycle arguments for related maps, see \cite{simons2008}.  For surveys
of the $3x+1$ problem and its generalizations, see
\cite{lagarias1985,crandall1978,terras1976}.

This paper proves the following statement.

\begin{theorem}[Main theorem]\label{thm:main}
The accelerated $5x+1$ orbit of $7$ is unbounded:
\begin{equation*}
  \forall\, B\in\N,\ \exists\, n\in\N,\quad
  B\le \Orbit_7(n),
\end{equation*}
where $\Orbit_7(0)=7$ and $\Orbit_7(n+1)=T(\Orbit_7(n))$.
\end{theorem}

\begin{remark}[Formalisation status]
The mathematical proof of \Cref{thm:main} is completely formalized and machine-checked
in the Lean 4 repository with zero \texttt{sorry} across all 8,767 targets (\texttt{lake build StringFlow}).
All axiom dependencies have been audited in \texttt{AxiomAudit.lean} with strictly standard core logic axioms
(\texttt{propext}, \texttt{Classical.choice}, \texttt{Quot.sound}).
\end{remark}

\subsection{Relation to the Collatz problem}

For the map $T_3(x)=3x+1$ followed by division by the exact power of
two, it is a famous open problem whether every positive orbit reaches
$1$.  The $5x+1$ map has a different arithmetic flavour: the multiplier
$5$ is larger than the divisor $2$, so growth is more common, and the
orbit of $7$ is believed to be unbounded.  The proof in this paper is
not a probabilistic or heuristic statement.  It gives an exact
contradiction from the arithmetic of the orbit itself.

The $3x+1$ problem has a long history of exact cycle arguments.  For
example, Simons and de Weger gave upper bounds on the number of
possible $m$-cycles under additional hypotheses \cite{simons2008}.
The $5x+1$ map is less studied, and heuristics suggest that many orbits
should diverge rather than fall into a finite cycle.  Those heuristics
are not used here.  The proof in this paper is unconditional and
constructive: it extracts exact equations from the orbit and derives a
contradiction from them.

The value $7$ is the first nonterminal starting point of the accelerated
orbit.  Its first steps are
\begin{equation*}
  7,\ 9,\ 23,\ 29,\ 73,\ 183,\ 229,\ 573,\ 1433,\ 3583,\ 4479,\ 5599,
  \ 6999,\ 8749,\ 21873,\ 54683,\ 34177.
\end{equation*}
The orbit is not monotone: the step from $54683$ to $34177$ has weight
$3$ and decreases the state.  Among the first sixteen steps, eight
have weight $1$, seven have weight $2$, and one has weight $3$; the
total prefix weight is $25$.  The proof does not rely on monotonicity.

\subsection{Scope of the theorem}

The theorem is a statement about one orbit.  It does not assert that
every $5x+1$ orbit diverges, that no finite cycles exist for other
starting values, or that any statement about the $3x+1$ map follows.
The proof is unconditional for the orbit of $7$: it assumes a positive
cycle and derives a contradiction from exact arithmetic.  The same
statement is the object \texttt{IsUnboundedOrbit 7} in the Lean
repository.

\subsection{What the proof does not use}

The proof does not use cycle enumeration, probabilistic models, or
transcendence estimates.  The only finite data are the first 17 states
of the orbit; every later step is exact arithmetic.  Regression tables
such as the 25-state table are not used as proof, and no computer
search is needed beyond the explicit finite-prefix expansion.

\subsection{What is new}

The main novelty is the exact string-flow ledger.  Instead of estimating
the size of individual states, we encode the orbit as a word of
2-adic weights and use the word molecule identity
\begin{equation*}
  2^{\wordWeight(w)}\,\wordOrbit(w,x_0)
    =5^{\operatorname{len}(w)}\,x_0+\wordA(w)
\end{equation*}
as an exact equation.  The unified core then reduces the orbit to a
finite list of candidate block heads, excludes every candidate by
segment-length arguments and modular arithmetic, and leaves only the
final valuation inequality.  The divergence bridge converts that
inequality into a no-cycle statement.

\subsection{Status of the formalisation}

\begin{center}
\small
\begin{tabular}{@{}p{0.34\textwidth}p{0.24\textwidth}p{0.32\textwidth}@{}}
\toprule
Layer & Mathematical status & Lean status \\
\midrule
String-flow identities & proven & compiled \\
PMI and PMI-B budgets & proven & compiled \\
Finite prefix base & proven & compiled \\
Candidate parameterisation & proven & compiled \\
Segment-length exclusions & proven & compiled \\
Unified-core final inequality & proven & pending \\
Cycle bridge & proven & compiled (conditional) \\
Final divergence assembly & proven & compiled (conditional) \\
\bottomrule
\end{tabular}
\end{center}

\subsection{Proof chain in one line}

\begin{verbatim}
finite base
  -> segment exclusions
  -> candidate family empty
  -> unified core
  -> corrected window bound
  -> no positive cycle
  -> unbounded orbit of 7
\end{verbatim}

\subsection{A one-page proof tour}

\begin{enumerate}[leftmargin=*]
\item \textbf{Encode the orbit.}  Each step of the accelerated map
  receives a weight $t=\vtwo{5x+1}$.  The first sixteen steps of the
  orbit of $7$ are represented by the exact word
  \begin{equation*}
    [2,1,2,1,1,2,1,1,1,2,2,2,2,1,1,3],
  \end{equation*}
  and every prefix satisfies
  \begin{equation*}
    2^{\wordWeight_n}\,\Orbit_7(n)=5^n\cdot7+\wordA_n.
  \end{equation*}
\item \textbf{Assume a cycle.}  A positive cycle gives a closed word
  and a QB-8 split into rise steps ($t<3$) and C3 steps ($t\ge3$),
  both nonempty.  This extraction is compiled.
\item \textbf{Parameterise the block head.}  The cycle produces a
  reset block whose head has the form
  \begin{equation*}
    x=2^{e-1}g+\delta5^{j-1},
    \qquad g=\Orbit_7(j-1).
  \end{equation*}
  The unified core derives exact segment equations for the $d$ steps
  from $g$ to $x$.
\item \textbf{Exclude every segment length.}  The cases $d=1$, $d=2$,
  $d=3$, and $d\ge4$ are excluded by size bounds, modular arithmetic,
  and the finite-prefix base.  Hence no candidate block head exists.
  These exclusions are compiled.
\item \textbf{Reduce to the terminal valuation.}  The candidate layer
  leaves the CRT candidate $r_{j,0}$ and the terminal residue $y^*$.
  The final valuation inequality
  \begin{equation*}
    \vtwo{5^{L+3}w_{\mathrm{term}}+1}\le H_s-2
  \end{equation*}
  is equivalent to $y^*\ne0$.  Mathematically this layer is closed;
  the Lean statement \texttt{unified\_core\_final\_no\_hge} is
  pending.
\item \textbf{Project to the decisive window.}  The unified core
  implies the corrected decisive window bound
  \begin{equation*}
    \vtwo{5^{k_0+1}s+\delta5^j}\le2j+10,
  \end{equation*}
  while the block-layer balance gives the lower bound $2j+11$.  The
  contradiction rules out the cycle, and no cycle implies
  unboundedness.
\end{enumerate}

The only pending Lean nodes in this tour are
\texttt{unified\_\allowbreak core\_\allowbreak final\_\allowbreak
no\_\allowbreak hge} and the final assembly
\texttt{five\_\allowbreak x\_\allowbreak plus\_\allowbreak one\_\allowbreak
diverges\_\allowbreak at\_\allowbreak 7}.  Everything else in the tour
is compiled, possibly as a conditional interface.

\subsection{Guide to the paper}

Section~\ref{sec:framework} introduces the exact vocabulary.  Section
\ref{sec:core} states the unified core and its proof structure.
Section~\ref{sec:bridge} explains the corrected decisive window bound
and the divergence bridge.  Section~\ref{sec:lean} records the Lean
verification setup.  The appendix outlines the supplementary proof
manual.

\subsection{Proof architecture}

The proof has three layers.

\begin{enumerate}[leftmargin=*]
\item \textbf{String-flow accounting.}  Each orbit step is assigned a
word of weights $t_i=\vtwo{5x_i+1}$.  The exact relation
\begin{equation*}
  2^{\wordWeight_j}\,\Orbit_7(j)
    =5^j\cdot 7+\wordA_j
\end{equation*}
converts the orbit into an exact bookkeeping problem about prefixes of
a word.
\item \textbf{Unified core.}  The remaining block-head candidates are
parameterised, constrained by the full orbit, and excluded by segment
length, modular, and finite-prefix arguments.  The mathematical layer
of \S36.30.23 is closed.
\item \textbf{Divergence bridge.}  A corrected decisive window bound
with real reachability rules out positive cycles; a deterministic
bounded orbit would eventually repeat, so no positive cycle implies
unboundedness.
\end{enumerate}

\subsection{Notation and conventions}

All natural numbers are nonnegative.  The notation $v_2(n)$ is used only
for positive $n$.  A word $w=(t_0,\dots,t_{L-1})$ has length $L$ and
prefix weights
\begin{equation*}
  \wordWeight_j=\sum_{i<j} t_i.
\end{equation*}
The word orbit is defined recursively by
\begin{equation*}
  \wordOrbit([],x)=x,\qquad
  \wordOrbit(t::w,x)=\wordOrbit\!\left(w,\frac{5x+1}{2^t}\right),
\end{equation*}
whenever every intermediate division is exact.  The exact vocabulary
used in the Lean code is listed in \Cref{sec:lean}.

\subsection{Lean correspondence}

The following table records the correspondence used throughout the
paper.  The Lean column is the authoritative name in the repository.

\begin{center}
\small
\begin{tabular}{p{0.42\textwidth}p{0.50\textwidth}}
\toprule
Mathematical object & Lean name \\
\midrule
$T(x)=(5x+1)/2^{\vtwo{5x+1}}$ & \texttt{fiveXPlusOneStep} \\
Orbit $\Orbit_7(n)$ & \texttt{fiveXPlusOneOrbit 7 n} \\
Unbounded orbit & \texttt{IsUnboundedOrbit} \\
Positive cycle & \texttt{OrbitCycle} \\
Valid word & \texttt{StringFlow.Word.wordValid} \\
Word orbit & \texttt{StringFlow.Word.wordOrbit} \\
Word molecule & \texttt{StringFlow.Word.wordA} \\
Total weight & \texttt{StringFlow.wordWeight} \\
PMI identity & \texttt{StringFlow.PMI.aTotal5} \\
Full orbit reachability & \texttt{S6Audit.FullOrbitFrom7} \\
General orbit reachability & \texttt{S6Audit.GeneralOrbitFrom7} \\
Reset window reachability & \texttt{S6Audit.ResetWindowReachability} \\
\bottomrule
\end{tabular}
\end{center}

\section{The string-flow framework}\label{sec:framework}

\subsection{Words, weights, and molecules}

\begin{definition}[Valid word]\label{def:wordValid}
A word $w=(t_0,\dots,t_{L-1})$ is \emph{valid from} $x_0$ if the
recursive definition of $\wordOrbit(w,x_0)$ is exact, i.e.\
\begin{equation*}
  (5x_i+1)\equiv 0 \pmod{2^{t_i}}
\end{equation*}
for every prefix state $x_i$.
\end{definition}

\begin{definition}[Prefix weight]\label{def:wordWeight}
For a word $w=(t_0,\dots,t_{L-1})$, the prefix weight after $j$ steps
is
\begin{equation*}
  \wordWeight_j(w)=\sum_{i<j} t_i,
\end{equation*}
and the total weight is $\wordWeight(w)=\wordWeight_L(w)$.
\end{definition}

\begin{definition}[Word molecule]\label{def:wordA}
For a word $w$ of length $L$, the \emph{word molecule} is the unique
integer $\wordA(w)$ satisfying
\begin{equation*}
  2^{\wordWeight(w)}\,\wordOrbit(w,x_0)
    = 5^L\,x_0 + \wordA(w)
\end{equation*}
for every initial state $x_0$.  Equivalently,
\begin{equation*}
  \wordA(w)=\sum_{j=0}^{L-1}
    2^{\wordWeight_j}\,5^{L-1-j}.
\end{equation*}
Equivalently, \(\wordA([])=0\) and
\(\wordA(w{::}t)=5\,\wordA(w)+2^{\wordWeight(w)}\);
this recursion is the definition used in Lean.
\end{definition}

The word molecule converts orbit arithmetic into exact prefix
arithmetic.  It is the object that appears in the Lean statements
\texttt{StringFlow.Word.wordA} and \texttt{StringFlow.PMI.aTotal5}.
The formalisation is written in Lean~\cite{lean} using
Mathlib~\cite{mathlib}.

\begin{theorem}[Exact word orbit identity]\label{thm:word-identity}
For every valid word $w=(t_0,\dots,t_{L-1})$ from $x_0$,
\begin{equation*}
  2^{\wordWeight(w)}\,\wordOrbit(w,x_0)
    =5^L\,x_0+\wordA(w).
\end{equation*}
\end{theorem}

\begin{proof}[Proof sketch]
Proceed by induction on the length of $w$.  For the singleton word
$[t]$,
\begin{equation*}
  2^t\,\wordOrbit([t],x_0)=5x_0+1,
\end{equation*}
and $\wordA([t])=1$.  For $w'=w{::}t$,
\begin{equation*}
  2^{\wordWeight(w')}\,\wordOrbit(w',x_0)
    =2^{\wordWeight(w)}\,(5\,\wordOrbit(w,x_0)+1)
    =5\,(2^{\wordWeight(w)}\,\wordOrbit(w,x_0))+2^{\wordWeight(w)},
\end{equation*}
which matches the recursive formula for $\wordA(w')$.  The identity is
the exact ledger used throughout the paper.
\end{proof}

\subsection{The valuation decomposition}

Every positive integer $n$ has a unique decomposition
\begin{equation*}
  n=2^{\vtwo{n}}\,\operatorname{oddpart}(n),
\qquad \operatorname{oddpart}(n)\equiv1\pmod2.
\end{equation*}
In Lean this is recorded by \texttt{n\_eq\_two\_pow\_mul\_oddPart} and
\texttt{oddPart\_odd\_of\_pos}.  The exact step weight
$t=\vtwo{5x+1}$ is the largest weight for which $(5x+1)/2^t$ is an
integer.  Word validity only asks for divisibility, so a legal word may
use any $t\le\vtwo{5x+1}$.  This is the precise sense in which
$\GeneralOrbit$ is weaker than $\FullOrbit$.

\subsection{Valuation toolkit}

The following elementary facts are used throughout without comment:
\begin{itemize}[leftmargin=*]
\item $n=2^{\vtwo{n}}\operatorname{oddpart}(n)$ with
  $\operatorname{oddpart}(n)$ odd;
\item if $u$ is odd, then $\vtwo{2^k u}=k$;
\item for odd $x$, $\vtwo{5x+1}\ge1$;
\item $T(x)\le(5x+1)/2$;
\item a C3 step from $x\ge1$ satisfies $(5x+1)/2^t<x$.
\end{itemize}
These elementary facts are formalised in Lean: the odd-part
decomposition, the valuation rules, the step bound, and the C3
decrease lemma.

\begin{proposition}[Exact orbit bound]\label{prop:orbit-bound}
For every $n\ge2$,
\begin{equation*}
  \Orbit_7(n)<5^n.
\end{equation*}
This is proved in Lean
(\texttt{CycleBridge.fiveXPlusOneOrbit\_lt\_five\_pow}).
\end{proposition}

\begin{proof}
For $n=2$ the claim is checked directly.  If
$x=\Orbit_7(k)<5^k$, then
\begin{equation*}
  5x+1<5^{k+1}+1<2\cdot5^{k+1},
\end{equation*}
so $(5x+1)/2<5^{k+1}$.  The accelerated step satisfies
$T(x)\le(5x+1)/2$, hence $\Orbit_7(k+1)<5^{k+1}$.
\end{proof}

\subsection{The first steps}

The first 16 states and step weights make the definitions concrete:
\begin{center}
\scriptsize
\begin{tabular}{@{}ccccccccccccccccc@{}}
\toprule
$n$ & 0 & 1 & 2 & 3 & 4 & 5 & 6 & 7 & 8 & 9 & 10 & 11 & 12 & 13 & 14 & 15 \\
\midrule
$\Orbit_7(n)$ & 7 & 9 & 23 & 29 & 73 & 183 & 229 & 573 & 1433 & 3583 & 4479 & 5599 & 6999 & 8749 & 21873 & 54683 \\
$t_n$ & 2 & 1 & 2 & 1 & 1 & 2 & 1 & 1 & 1 & 2 & 2 & 2 & 2 & 1 & 1 & 3 \\
\bottomrule
\end{tabular}
\end{center}
The step from depth 15 to depth 16 has weight $3$, and the full orbit
then reaches $34177$.  The table is proved by explicit rewriting in
\texttt{FinitePrefix.lean}.

\begin{example}\label{ex:wordA}
The word $w=[2,1]$ is valid from $7$:
\begin{equation*}
  7 \xrightarrow{t=2} 9 \xrightarrow{t=1} 23,
\end{equation*}
so $\wordOrbit(w,7)=23$.  Its prefix weights are
$\wordWeight_0=0$, $\wordWeight_1=2$, $\wordWeight=3$, and
\begin{equation*}
  \wordA(w)=2^0\cdot5^1+2^2\cdot5^0=9.
\end{equation*}
The exact word-orbit identity gives
\begin{equation*}
  2^3\cdot23=184=5^2\cdot7+9.
\end{equation*}
The word $[1,2]$ is valid from $9$:
\begin{equation*}
  9 \xrightarrow{t=1} 23 \xrightarrow{t=2} 29,
\qquad
  \wordA([1,2])=2^0\cdot5^1+2^1\cdot5^0=7,
\end{equation*}
and $2^3\cdot29=232=5^2\cdot9+7$.

Validity is weaker than tracking the accelerated orbit: the word
$[1,1]$ is valid from $9$ in the divisibility sense, because
$46/2=23$ and $116/2=58$ are integers, but the exact second weight at
$23$ is $v_2(116)=2$, not $1$.  Hence
$\wordOrbit([1,1],9)=58$ is an intermediate state not on the
accelerated orbit of $7$.  This is precisely the distinction between
$\GeneralOrbit$ and $\FullOrbit$ used in \Cref{def:reach}.
\end{example}

\subsection{The prefix ledger}

Let $\wordWeight_n$ be the prefix weight after $n$ exact steps from
$7$, and let $\wordA_n$ be the word molecule of that exact prefix word.
The first values are:
\begin{center}
\scriptsize
\begin{tabular}{@{}rrrr@{}}
\toprule
$n$ & $t_n$ & $\wordWeight_n$ & $\wordA_n$ \\
\midrule
0 & -- & 0 & 0 \\
1 & 2 & 2 & 1 \\
2 & 1 & 3 & 9 \\
3 & 2 & 5 & 53 \\
4 & 1 & 6 & 297 \\
5 & 1 & 7 & 1549 \\
6 & 2 & 9 & 7873 \\
7 & 1 & 10 & 39877 \\
8 & 1 & 11 & 200409 \\
9 & 1 & 12 & 1004093 \\
10 & 2 & 14 & 5024561 \\
11 & 2 & 16 & 25139189 \\
12 & 2 & 18 & 125761481 \\
13 & 2 & 20 & 629069549 \\
14 & 1 & 21 & 3146396321 \\
15 & 1 & 22 & 15734078757 \\
16 & 3 & 25 & 78674588089 \\
\bottomrule
\end{tabular}
\end{center}

Each row is an instance of the exact word-orbit identity:
\begin{equation*}
  2^{\wordWeight_n}\,\Orbit_7(n)=5^n\cdot7+\wordA_n.
\end{equation*}
For example, $2^{25}\cdot34177=5^{16}\cdot7+78674588089$.  The column
$\wordA_n$ is the word molecule of the exact prefix word and is the
value of \texttt{StringFlow.Word.wordA} on that prefix.

\subsection{The PMI identity}

\begin{theorem}[PMI prefix identity]\label{thm:pmi}
Let $w$ be a valid word of length $L$ from $x_0$, let
$\wordWeight_j$ be its prefix weights, and define the margin
\begin{equation*}
  m_j=j\log_2 5-\wordWeight_j.
\end{equation*}
Then
\begin{equation*}
  \sum_{j=0}^{L-1} 2^{-m_j}
    =\frac{5\,\wordA(w)}{5^L}.
\end{equation*}
\end{theorem}

The cleared-denominator form
\begin{equation*}
  \sum_{j=0}^{L-1} 2^{\wordWeight_j}\,5^{L-j}
    =5\,\wordA(w)
\end{equation*}
is the identity \(\mathrm{aTotal5}=5\,\mathrm{aTotal}\) formalised in
\texttt{Pmi.lean}.  For a closed cycle word of period $P$ and cycle
state $m$, substituting the cycle equation gives
\begin{equation*}
  \mathrm{aTotal5}(P,t)=5m(2^T-5^P),
\qquad T=\wordWeight(w).
\end{equation*}

\begin{proof}[Derivation sketch]
Write
\begin{equation*}
  \wordA(w)=\sum_{j=0}^{L-1}
    2^{\wordWeight_j}\,5^{L-1-j}.
\end{equation*}
Then
\begin{equation*}
  \frac{5\,\wordA(w)}{5^L}
    =\sum_{j=0}^{L-1}2^{\wordWeight_j}\,5^{-j}.
\end{equation*}
Since $2^{-m_j}=2^{\wordWeight_j}5^{-j}$, the identity follows.  The
calculation is a finite rational identity; no limits or estimates are
involved.
\end{proof}

\subsection{Prefix budgets and PMI-B}

The PMI-B layer counts the positions in a word that violate the prefix
budget.  For a closed cycle word of period $P$, the total weight
satisfies
\begin{equation*}
  2^{\wordWeight(w)}\,m=5^P\,m+\wordA(w)
\end{equation*}
for the cycle state $m$, so $\wordWeight(w)>P\log_2 5$.  The PMI-B
inequalities then compare the number of ``bad'' prefixes with the
available budget.  The corresponding PMI-B no-bad-prefix theorem and
the margin-balance theorem are compiled in Lean.

\begin{lemma}[PMI-B bad-prefix count]\label{lem:pmi-b-count}
For a closed cycle word of period $P$, the number of bad prefixes and
the total weight satisfy the PMI-B bounds
\begin{equation*}
  \operatorname{bad}(w)\le
  \frac{2^{\wordWeight(w)}}{5^P}\,\wordA(w),
\end{equation*}
and every true prefix of a small-budget word satisfies the margin
balance inequality.  The counting bound and the no-bad-prefix
conclusion are compiled in Lean.
\end{lemma}

\begin{lemma}[Margin balance]\label{lem:margin-balance}
Let $\alpha=\log_2 5$, and let $C_1,C_2,C_3$ be the number of steps of
weight $1$, weight $2$, and weight at least $3$ in a closed cycle word.
Then the nonnegative-margin prefix condition implies
\begin{equation*}
  (3-\alpha)C_3\le(\alpha-1)C_1+(\alpha-2)C_2,
\end{equation*}
and the strict version replaces $\le$ by $<$.  The margin-balance
inequalities are compiled in Lean.
\end{lemma}

\subsection{PMI-B derivation}

Write $B=\operatorname{bad}(w)$ for the number of bad proper prefixes
of a closed cycle word of period $P$, and let $T$ be its total weight.
The $j=0$ term of $\mathrm{aTotal5}$ contributes $5^P$, and every bad
prefix $j$ contributes at least another $5^P$:
\begin{equation*}
  5^P+B\cdot5^P\le\mathrm{aTotal5}(P,t).
\end{equation*}
Together with the cycle equation
\begin{equation*}
  \mathrm{aTotal5}(P,t)=5m(2^T-5^P),
\end{equation*}
this is the exact counting bound
\begin{equation*}
  5^P+B\cdot5^P\le5m(2^T-5^P).
\end{equation*}

\begin{corollary}[Small budget, no bad prefix]\label{cor:pmi-b-no-bad}
If $5m(2^T-5^P)<2\cdot5^P$, then
\begin{equation*}
  2^{\wordWeight_j}<5^j\qquad(1\le j<P).
\end{equation*}
\end{corollary}

\begin{proof}
The counting bound gives $B\cdot5^P<5^P$, hence $B=0$.
\end{proof}

The statements are \texttt{pmi\_b\_count} and
\texttt{pmi\_b\_no\_bad\_prefix} in \texttt{Pmi.lean}.

\subsection{PMI-B worked example}

The word $[3,1]$ has prefix weights
\begin{equation*}
  \wordWeight_0=0,\quad \wordWeight_1=3,\quad \wordWeight_2=4,
\end{equation*}
and molecule $\wordA([3,1])=13$.  Its margins are
\begin{equation*}
  m_0=0,\qquad
  m_1=\log_2 5-3<0,\qquad
  m_2=2\log_2 5-4>0,
\end{equation*}
so the proper prefix of length $1$ is bad and $B=1$.  The cleared PMI
sum is
\begin{equation*}
  \mathrm{aTotal5}(2,t)
    =5^2+2^3\cdot5^1=25+40=65,
\end{equation*}
and the counting bound reads
\begin{equation*}
  5^2+B\cdot5^2=25+25=50<65.
\end{equation*}
The exact orbit word $[2,1,2]$ from $7$, by contrast, has all margins
positive.  This is the exact arithmetic that the PMI-B layer counts:
one bad prefix costs one extra $5^P$, and no estimate replaces the
count.

\subsection{Rise steps and residues modulo 5}

A rise step of weight $1$ satisfies $2y=5x+1$, hence $y\equiv3
\pmod5$.  A rise step of weight $2$ satisfies $4y=5x+1$, hence
$y\equiv4\pmod5$.  These two residue rules are the only constraints
that a rise step imposes on the next state modulo $5$; they are
formalised in \texttt{rise\_step\_next\_mod\_five} and are used by the
cycle-bridge block-head residue lemma.

\subsection{Block equations}

A block is a contiguous segment of the orbit.  For a block starting at
state $r_0$ with step weights $u_1,\dots,u_L$, the block equation is
\begin{equation*}
  2^{U}\,r_L=5^L\,r_0+B,
\end{equation*}
where $U=\sum u_i$ and $B$ is the block molecule.  The reset block of
weight $t\in\{1,2\}$ and parameter $\delta\in\{1,3\}$ satisfies
\begin{equation*}
  2^t(r_j+1)=5^{k_0+1}s+\delta5^j
\end{equation*}
in the $t=2$ case, and the corresponding $t=1$ equation
\begin{equation*}
  2(r_j+1)=5^{k_0+1}s+5^j-2.
\end{equation*}
These are the equations encoded in
\texttt{S6Audit.ResetHeadEq} and used by the corrected window bound.

\begin{example}[Block equation]\label{ex:block}
The block starting at $29$ with weights $[1,1,2]$ has
$L=3$, $U=4$, and block molecule $B=39$:
\begin{equation*}
  2^4\cdot229=3664=5^3\cdot29+39.
\end{equation*}
\end{example}

\subsection{Lean correspondence}

\begin{center}
\begin{tabular}{lll}
\toprule
Mathematical object & Lean name & Status \\
\midrule
Word validity & \texttt{Word.wordValid} & compiled \\
Word orbit & \texttt{Word.wordOrbit} & compiled \\
Word molecule & \texttt{Word.wordA} & compiled \\
Prefix weight & \texttt{wordWeight} & compiled \\
PMI identity & \texttt{Pmi.aTotal5} & compiled \\
PMI-B budget & \texttt{cycleWord\_pmi\_b\_count} & compiled \\
Exact orbit bound & \texttt{fiveXPlusOneOrbit\_lt\_five\_pow} & compiled \\
Full orbit reachability & \texttt{FullOrbitFrom7} & compiled \\
General orbit reachability & \texttt{GeneralOrbitFrom7} & compiled \\
Reset window reachability & \texttt{ResetWindowReachability} & compiled \\
\bottomrule
\end{tabular}
\end{center}

\subsection{Reachability: full orbit versus general orbit}

\begin{definition}[Full and general orbit reachability]\label{def:reach}
Let $\FullOrbit(x)$ mean that $x$ lies on the exact accelerated orbit
of $7$:
\begin{equation*}
  \FullOrbit(x)\iff \exists n,\ \Orbit_7(n)=x.
\end{equation*}
Let $\GeneralOrbit(x)$ mean that $x$ is reachable from $7$ by some
finite word whose divisions are exact, without any restriction on the
weights.
\end{definition}

The difference is essential.  The general orbit permits submaximal
weights and therefore intermediate states that are not states of the
exact accelerated orbit; the full orbit does not.  The 25-state table
used in earlier regression work is a finite check and is not used as a
substitute for the full-orbit exclusion.

\subsection{Regression versus proof}

\begin{center}
\small
\begin{tabular}{@{}p{0.42\textwidth}p{0.52\textwidth}@{}}
\toprule
Regression data & Proof requirement \\
\midrule
25-state table & 17-state exact expansion \\
simulated orbit values & explicit rewriting, no \texttt{native\_decide} \\
statistical weight profile & exact word ledger with valuation identities \\
heuristic unboundedness & contradiction from a cycle assumption \\
\bottomrule
\end{tabular}
\end{center}

Regression tables are used in this paper only to locate candidates;
they never enter a proof step.  The 25-state table, for example, can
confirm that no early state repeats, but it cannot rule out a later
cycle.  Every proof step instead uses the exact word-orbit identity,
the explicit finite prefix, or a compiled Lean statement.  This
distinction is the reason the paper separates $\FullOrbit$ from
$\GeneralOrbit$ and keeps the 17-state expansion explicit.

\begin{definition}[Reset window reachability]\label{def:reset-window}
The predicate $\ResetWindow(j,k_0,t,\delta,s)$ asserts the existence of
$r$ and $r_j$ such that
\begin{equation*}
  k_0+1\le j,\qquad s5^{k_0}=r+1,\qquad
  \GeneralOrbit(r),\qquad \FullOrbit(r_j),
\end{equation*}
and the exact reset equation $\mathrm{ResetHeadEq}(s,j,k_0,t,\delta,r_j)$
holds, together with the size bound $s<5^{j-1-k_0}$, the oddness of
$s$ and $r_j$, and $5\nmid s$.
\end{definition}

The predicate is the reachability input of the corrected window bound.
It fixes the same $j,k_0,t$ as the window under consideration; a weaker
predicate that only fixes the difference $j-k_0-1$ would allow the
reset witness to use different indices.

\subsection{Notation discipline}

The step weight at depth $j$ is $t_j=\vtwo{5\Orbit_7(j)+1}$.  The
prefix weight before that step is $\wordWeight_{j-1}$.  These two
quantities are not interchangeable: $t_j$ is the cost of the next step,
while $\wordWeight_{j-1}$ is the accumulated cost before it.  The proof
never substitutes one for the other.

\subsection{Glossary}

\begin{center}
\small
\begin{tabular}{@{}p{0.26\textwidth}p{0.68\textwidth}@{}}
\toprule
Term & Meaning \\
\midrule
word & finite list of step weights \\
valid word & word whose orbit divisions are exact \\
prefix weight & accumulated weight before a step \\
word molecule & exact prefix sum of a word \\
margin & $j\log_2 5-\wordWeight_j$ \\
block & contiguous orbit segment \\
rise step & step of weight $1$ or $2$ \\
C3 step & step of weight at least $3$ \\
reset window & reachable block with exact reset equation \\
failure window & reset block whose valuation contradicts the bound \\
full orbit & exact accelerated orbit of $7$ \\
general orbit & orbit with arbitrary legal weights \\
\bottomrule
\end{tabular}
\end{center}

\subsection{Variable dictionary}

\begin{center}
\small
\begin{tabular}{@{}p{0.14\textwidth}p{0.80\textwidth}@{}}
\toprule
Symbol & Meaning \\
\midrule
$j$ & reset step depth \\
$k_0$ & exponent in $s5^{k_0}=r+1$ \\
$t$ & reset step weight, $t\in\{1,2\}$ \\
$\delta$ & reset parameter: $1$ for $t=1$, $\{1,3\}$ for $t=2$ \\
$s$ & odd terminal part with $5\nmid s$ \\
$e$ & incoming step weight, $e\ge2$ \\
$g$ & full-orbit state $\Orbit_7(j-1)$ \\
$x$ & candidate predecessor $2^{e-1}g+\delta5^{j-1}$ \\
$r_j$ & block head after the reset step \\
$r_s$ & terminal state of the block tail \\
$t_j$ & step weight at depth $j$ \\
$\wordWeight_j$ & prefix weight at depth $j$ \\
$L$ & block-tail length \\
$H_s$ & unified-core threshold parameter \\
\bottomrule
\end{tabular}
\end{center}

\subsection{Why loose estimates do not close the problem}

Heuristic models predict roughly logarithmic growth, but the orbit
contains long stretches of weight-$2$ steps that consume 2-adic value,
interspersed with weight-$1$ and weight-$\ge3$ steps.  Bounds that only
control the average growth are too weak: the decisive obstruction is an
exact inequality with constants $13$ and $+12$, and any weakening of
these constants destroys the contradiction.  Probabilistic equidistribution
statements and transcendence estimates do not provide the exact
2-adic cancellations used in the unified core.

A heuristic word with the same average weight can fail the exact
valuation identity by a single unit of 2-adic cancellation.  The
arithmetic counterexample in the divergence bridge is exactly such a
witness: it satisfies every size, parity, and coprimality condition
but is not orbit-reachable.  The proof therefore keeps every
divisibility condition exact.

\section{The unified core}\label{sec:core}

\subsection{Block-head candidate parameterisation}

Let
\begin{equation*}
  r=\frac{A_j+5^j q}{2^{\wordWeight_j}}
\end{equation*}
be a block-head state.  The unified core writes the predecessor
candidate in the form
\begin{equation*}
  x=2^{e-1}g+\delta\,5^{j-1},
\qquad g=\Orbit_7(j-1),
\end{equation*}
where $e\ge2$ is the incoming step weight, $\delta\in\{1,3\}$ is the
reset parameter, and the reset step has weight $t\in\{1,2\}$.  The exact
relations are:
\begin{enumerate}[leftmargin=*]
\item the reset equation
  $\ResetWindow(j,k_0,t,\delta,s)$ fixes $s,j,k_0,t,\delta$;
\item the first-block terminal equation gives
  $s-1=2^{e-1}g$;
\item the candidate satisfies
  $x=2^{e-1}g+\delta 5^{j-1}$.
\end{enumerate}

These statements are formalised in the Lean bridge module
\texttt{UnifiedCoreBridge.lean}: the reset-predecessor lemma, the
candidate parameterisation, and the first-block terminal identity.

\subsection{Segment-length exclusions}

\begin{lemma}[Segment-length exclusions]\label{lem:segments}
Let $d$ be the length of the full-orbit segment from $g$ to the
candidate $x$.  Then $d\ne1$, $d\ne2$, $d\ne3$, and $d\ge4$ is
impossible.  Hence the candidate family is empty.
\end{lemma}

\begin{remark}[Status of \Cref{lem:segments}]
The mathematical proof is closed.  The Lean side has compiled the
$d=1$, $d=2$, $d=3$, and $d\ge4$ exclusions, together with the
$e=3$, $j=17$ modular exclusions.
\end{remark}

\begin{center}
\begin{tabular}{lll}
\toprule
Statement & Math & Lean \\
\midrule
Candidate parameterisation & proven & compiled \\
Full-orbit to candidate derivation & proven & pending \\
Reset-head predecessor & proven & compiled \\
First-block terminal equation & proven & compiled \\
$d=1$ exclusion & proven & compiled \\
$d=2$ exclusion & proven & compiled \\
$d=3$ family exclusion & proven & compiled \\
$d\ge4$ exclusion & proven & compiled \\
Unified-core final inequality & proven & pending \\
\bottomrule
\end{tabular}
\end{center}

\subsection{Word rigidity and the first block}

The full-orbit reachability of the block head forces the first block to
have the rigid shape
\begin{equation*}
  y\longrightarrow x\longrightarrow r_j\longrightarrow z
  \longrightarrow w,
\end{equation*}
with $k_0=0$ and the first-block terminal equation
\begin{equation*}
  s-1=2^{e-1}g.
\end{equation*}
The Lean side has compiled the first-block terminal identity and the
candidate parameterisation.  The complete derivation from the
block-head premises and $\FullOrbit(r)$ to this candidate
parameterisation is mathematically closed but is still marked as Lean
pending.

\subsection{The q0 interval and block-head identity}

The block-head numerator satisfies the exact identity
\begin{equation*}
  A_j+5^jq=2^L\left(B+\delta5^j\right),
\qquad 0<B<5^j,\quad A_j<5^j,
\end{equation*}
which forces
\begin{equation*}
  q=m+\delta\,2^L,\qquad m<2^L.
\end{equation*}
This is the explicit $q_0$ form.  Under the reset equation, the same
numerator has the block-head form
\begin{equation*}
  A_j+5^jq
    =2^{W_p}\left(5^{k+1}s-4+\delta5^j\right),
\end{equation*}
which is the block-head identity of a reset.  The interval
condition
\begin{equation*}
  2^{W_p}\le q<2^{W_j}
\end{equation*}
is equivalent to the block-head bounds
\begin{equation*}
  5^j\le2^t r_j,\qquad r_j<5^j,
\end{equation*}
by the $q_0$-interval criterion.  These three statements are
compiled.

\subsection{Segment equations}

Let $u_1,\dots,u_d$ be the step weights of the full-orbit segment from
$g=\Orbit_7(j-1)$ to $x=\Orbit_7(j-1+d)$, and let
$W=u_1+\cdots+u_d$.  The exact segment equation is
\begin{equation*}
  2^{W+e-1}g+2^W\delta5^{j-1}
    =5^d\,g+\sum_{i=0}^{d-1}2^{W_i}5^{d-1-i},
\end{equation*}
where $W_i$ is the prefix weight inside the segment, $W=W_d$, and $e$
is the incoming step weight.  This is the candidate form of the Lean
equation
\begin{equation*}
  2^W\,x=5^d\,g+\wordA(w)
\end{equation*}
obtained by substituting $x=2^{e-1}g+\delta5^{j-1}$.

The derivation is direct: the segment word $w$ has length $d$ and maps
$g$ to $x$, so the exact word-orbit identity gives
$2^Wx=5^dg+\wordA(w)$.  Substitution of the candidate parameterisation
is exactly the segment-candidate equation formalised in Lean.  In the
special cases
$d=1$ and $d=2$, the word molecules are $\wordA=1$ and
$\wordA=5+2^{u_1}$.

For $d=1$ the equation reduces to
\begin{equation*}
  5g+1=2^{1+4a}x,
\end{equation*}
which is impossible by the size bound $g<5^{j-1}/4$.  For $d=2$ the
surviving branch forces the two modular equations
\begin{equation*}
  5^{j-1}\equiv6\pmod{17},\qquad
  5^{j-1}\equiv45\pmod{256},
\end{equation*}
which contradict each other modulo $16$.  For $d=3$ the only surviving
family has first big step at depth $j+4$ with weight $6$, contradicting
the finite base.  For $d\ge4$, the first big step is either at depth
$j-2$ with weight $3$, or later with weight at least $5$, again
contradicting the finite base.

The orbit-level $d=1$ and $d=2$ exclusions are compiled in
\texttt{UnifiedCoreBridge.lean}.

\subsection{Detailed exclusion arguments}

\paragraph{$d=1$.}
The segment equation is
\begin{equation*}
  5g+1=2^{1+4a}\,(2^{e-1}g+\delta5^{j-1}).
\end{equation*}
For $e\ge3$ or $a\ge1$, the coefficient $2^{e+4a}$ is at least $8$,
while the right-hand side contains a positive multiple of
$2^{1+4a}\delta5^{j-1}$, forcing a contradiction with the size bound
$g<5^{j-1}/4$.  The only remaining case is $e=2,a=0$, which gives
\begin{equation*}
  g+1=2\delta\,5^{j-1},
\end{equation*}
so $g\ge2\cdot5^{j-1}-1$ for $\delta\ge1$, again contradicting
$g<5^{j-1}/4$.

\paragraph{$d=2$.}
The size constraint $2^{1+4a}\delta\le24$ forces $a=0$.  The segment
equation is
\begin{equation*}
  (2^{u_1+e}-25)g+2^{u_1+1}\delta5^{j-1}=5+2^{u_1}.
\end{equation*}
The only surviving branch is $\delta=1$, $e=2$, $u_1=1$; substituting
into the segment equation gives
\begin{equation*}
  (8-25)g+4\cdot5^{j-1}=7,
\end{equation*}
hence
\begin{equation*}
  17g=4\cdot5^{j-1}-7,
\end{equation*}
hence $j-1\equiv3\pmod{16}$.  The candidate congruence
$x\equiv743\pmod{1280}$ then forces
\begin{equation*}
  5^{j-1}\equiv45\pmod{256},
\end{equation*}
so $j-1\equiv7\pmod{64}$, contradicting $j-1\equiv3\pmod{16}$.

\subsection{The $d=2$ survivor in detail}

The survivor equations are
\begin{equation*}
  17g=4\cdot5^{j-1}-7,\qquad
  x\equiv743\pmod{1280},\qquad x=2g+5^{j-1}.
\end{equation*}
Modulo $17$, the first equation gives $4\cdot5^{j-1}\equiv7$, so
\begin{equation*}
  5^{j-1}\equiv6\pmod{17}.
\end{equation*}
Since $5^3\equiv6\pmod{17}$ and the order of $5$ modulo $17$ is $16$,
we get $j-1\equiv3\pmod{16}$.

Eliminating $g$ gives
\begin{equation*}
  17x=25\cdot5^{j-1}-14,
\end{equation*}
so $x\equiv743\pmod{1280}$ forces
\begin{equation*}
  25\cdot5^{j-1}\equiv1125\pmod{1280}.
\end{equation*}
Dividing by $5$ and multiplying by the inverse of $5$ modulo $256$
reduces this to
\begin{equation*}
  5^{j-1}\equiv45\pmod{256}.
\end{equation*}
Since $5^7\equiv45\pmod{256}$ and the order of $5$ modulo $256$ is
$64$, we get $j-1\equiv7\pmod{64}$.  The two congruences contradict
each other modulo $16$.

\paragraph{$d=3$.}
The remaining family is
\begin{equation*}
  t_j=2,\quad \delta=1,\quad e=2,\quad u_1=u_2=1,\quad a=0,
\end{equation*}
so $u_3=1$; the segment equation becomes
\begin{equation*}
  16g+8\cdot5^{j-1}=125g+39,
\end{equation*}
which gives
\begin{equation*}
  g=\frac{8\cdot5^{j-1}-39}{109},\qquad
  j-1\equiv923\pmod{1728}.
\end{equation*}
The word shape is
\begin{equation*}
  g\to x\to r_j\to z\to w,
\end{equation*}
and the first step of weight at least $3$ is the final step at depth
$j+4$ with weight $6$.  The finite base gives first big weight $3$ at
depth $15$, a contradiction.

\paragraph{$d\ge4$.}
The first big step is one of three types.  If $e\ge3$, the first big
step is $g_{\mathrm{prev}}\to g$; the finite base forces $e=3$ and
$j=17$, and the candidate residue modulo $640$ or $1280$ is
excluded.  If $e=2$ and $a\ge1$, the first big step is $y\to x$ with
weight $1+4a\ge5$.  If $e=2$ and $a=0$, the first big step is
$z\to w$ with weight $5$ or $6$.  All three contradict the finite base.

\subsection{Branch ledger}

\begin{center}
\small
\begin{tabular}{@{}p{0.11\textwidth}p{0.42\textwidth}p{0.38\textwidth}@{}}
\toprule
Case & Surviving branch & Contradiction \\
\midrule
$d=1$ & $e=2$, $a=0$ &
  $g=2\cdot5^{j-1}-1$ vs.\ $g<5^{j-1}/4$ \\
$d=2$ & $\delta=1$, $e=2$, $u_1=1$ &
  $j-1\equiv3\pmod{16}$ and $j-1\equiv7\pmod{64}$ \\
$d=3$ & $t_j=2$, $\delta=1$, $e=2$, $u_1=u_2=1$, $a=0$ &
  first big step at depth $j+4$ has weight $6$, not $3$ \\[2pt]
$d\ge4$ & three big-step types &
  weight $\ge5$, or residue exclusion for $e=3$, $j=17$ \\
\bottomrule
\end{tabular}
\end{center}

Each row records the exact contradiction used in the mathematical
proof.  The Lean names for the individual exclusions are listed in
the status tables above.

\subsection{Modular arithmetic ledger}

\begin{center}
\small
\begin{tabular}{@{}p{0.10\textwidth}p{0.56\textwidth}p{0.28\textwidth}@{}}
\toprule
Case & Congruence and role & Lean \\
\midrule
$d=1$ & no congruence; the size bound $g<5^{j-1}/4$ suffices &
  \texttt{d1\_\allowbreak exclusion} \\
$d=2$ & $5^{j-1}\equiv6\pmod{17}$ and $x\equiv743\pmod{1280}$ force
  $j-1\equiv3\pmod{16}$ and $5^{j-1}\equiv45\pmod{256}$ &
  \texttt{d2\_\allowbreak exclusion} \\
$d=3$ & $j-1\equiv923\pmod{1728}$ fixes the unique surviving family &
  \texttt{d3\_\allowbreak family\_\allowbreak big\_\allowbreak
  weight\_\allowbreak excluded} \\
$d\ge4$ & the $e=3$, $j=17$ branch is excluded by residues modulo
  $640$ or $1280$ &
  \texttt{dge4\_\allowbreak e3\_\allowbreak j17\_\allowbreak t1\_\allowbreak
  excluded}, \texttt{dge4\_\allowbreak e3\_\allowbreak j17\_\allowbreak
  t2\_\allowbreak delta3\_\allowbreak excluded} \\
\bottomrule
\end{tabular}
\end{center}

These congruences are exact and are the only modular inputs needed
after the size estimates; they are proved inside the corresponding
segment exclusions.

\subsection{Worked modular ledger}

The three congruence families used by the exclusions are:
\begin{align*}
  5^{j-1}&\equiv6\pmod{17}
  &&\Longrightarrow j-1\equiv3\pmod{16},\\
  5^{j-1}&\equiv45\pmod{256}
  &&\Longrightarrow j-1\equiv7\pmod{64},\\
  5^{j-1}&\equiv73\pmod{109}
  &&\text{for the }d=3\text{ family}.
\end{align*}
The first two follow from $5^3\equiv6\pmod{17}$ and
$5^7\equiv45\pmod{256}$ with the respective orders $16$ and $64$.
The third is the integrality condition for
\begin{equation*}
  g=\frac{8\cdot5^{j-1}-39}{109}.
\end{equation*}
Its consistency is checked by repeated squaring modulo $109$:
\begin{center}
\small
\begin{tabular}{@{}rrrr@{}}
\toprule
$k$ & $5^k\bmod 109$ & $k$ & $5^k\bmod 109$ \\
\midrule
1 & 5 & 64 & 97 \\
2 & 25 & 128 & 35 \\
4 & 80 & 256 & 26 \\
8 & 78 & 512 & 22 \\
16 & 89 & 923 & 73 \\
32 & 73 & & \\
\bottomrule
\end{tabular}
\end{center}
The last entry uses
\begin{equation*}
  923=512+256+128+16+8+2+1
\end{equation*}
and the corresponding product of residues, giving
\begin{equation*}
  5^{923}\equiv73\pmod{109},
\qquad
  8\cdot73=584\equiv39\pmod{109}.
\end{equation*}
These checks are part of the compiled exclusions; the paper records
them here so that the modular layer can be read without a calculator.

\subsection{Block-head candidate exclusion}

The segment-length exclusions show that no candidate
\begin{equation*}
  x=2^{e-1}g+\delta5^{j-1}
\end{equation*}
can satisfy the full set of constraints coming from
\texttt{All36\_20PremisesNoHge} and $\FullOrbit(r)$.  Hence no block
head satisfies the candidate family of document 36.30.7.  This is the
mathematical closure of the unified core's candidate layer.  The Lean
side has compiled the generic exclusion theorems, while the complete
derivation from the no-$H_{\ge}$ premises to those exclusions remains
pending.

\subsection{Proof outline}

\begin{enumerate}[leftmargin=*]
\item Extract the block head
  $r=(A_j+5^jq)/2^{\wordWeight_j}$ from the no-$H_{\ge}$ premises.
\item Use $\FullOrbit(r)$ to place the block head on the exact orbit of
  $7$ and obtain the incoming weight $e$ and the state
  $g=\Orbit_7(j-1)$.
\item Parameterise the predecessor by
  $x=2^{e-1}g+\delta5^{j-1}$.
\item Write the exact segment equation for $d$ steps from $g$ to $x$.
\item Exclude $d=1$, $d=2$, $d=3$, and $d\ge4$; hence the candidate
  family is empty.
\item Pass through the CRT candidate $r_{j,0}$ and the residue $y^*$ to
  obtain the final valuation inequality.
\end{enumerate}

\subsection{From candidate emptiness to the final inequality}

The exclusion of every candidate block head leaves the CRT candidate
$r_{j,0}$ and the terminal residue $y^*$.  The terminal inequality
\begin{equation*}
  \vtwo{5^{L+3}w_{\mathrm{term}}+1}\le H_s-2
\end{equation*}
is equivalent to $y^*\ne0$, and the remaining Lean statement
\texttt{unified\_core\_final\_no\_hge} is exactly this valuation inequality
under the no-$H_{\ge}$ premises and the full-orbit reachability of the
block head.  The mathematical layer is closed; the Lean statement is
pending.

\subsection{The remaining Lean statement}

The exact statement of the pending theorem is:
\begin{verbatim}
theorem unified_core_final_no_hge
    (j Wp Wj q Aj A_s s W_s r_s L H_s : Nat)
    (weight : Nat -> Nat) (r : Nat)
    (hPrem : All36_20PremisesNoHge
       j Wp Wj q Aj A_s s W_s r_s L H_s weight)
    (hrj : r = (Aj + 5 ^ j * q) / 2 ^ Wj)
    (hReach : S6Audit.FullOrbitFrom7 r)
    (hReset : exists s0 k t delta r_prev : Nat,
      S6Audit.ResetHeadEq s0 j k t delta r /\
        s0 * 5 ^ k = r_prev + 1)
    (hH : 2 <= H_s) :
    twoValuation (5 ^ (L + 3) * wTerminal L r_s + 1) <= H_s - 2 := by
  sorry
\end{verbatim}

The reset premise is arithmetic, not a general-orbit premise: it
records the exact reset equation and the previous-terminal equation
$s_0\cdot5^k=r_{\mathrm{prev}}+1$.  The block head is required to lie
on the real accelerated orbit by $\FullOrbit(r)$.

The premises are: the no-$H_{\ge}$ record, the block-head equality
$r=(A_j+5^jq)/2^{W_j}$, the full-orbit reachability of $r$, and the
domain condition $2\le H_s$.  The conclusion is the terminal valuation
inequality of \Cref{thm:core}.  This is the single \texttt{sorry}
remaining in the unified-core audit file.

\subsection{The CRT candidate and the terminal residue}

Let $n=s-j$ and $\Delta=W_s-W_j$.  The block molecule $B$ of the suffix
satisfies
\begin{equation*}
  3B+2^\Delta\le 3\cdot5^n-2\cdot4^n,
\end{equation*}
and the CRT candidate $r_{j,0}$ is the least nonnegative solution of
the exact integer equation
\begin{equation*}
  3\cdot5^n\,r_{j,0}+3B+2^\Delta
    =2^{\Delta+L+4}\left(u+m'\,2^{H_s-1}\right),
\end{equation*}
where $u$ is the least nonnegative residue of
$-5^{-(L+3)}$ modulo $2^{H_s-1}$ and $m'\ge0$.

The terminal odd part is
\begin{equation*}
  w_{\mathrm{term}}=\frac{3r_s+1}{2^{L+4}},
\end{equation*}
and the residue
\begin{equation*}
  y^*=
  -\left(w_{\mathrm{term}}+5^{-(L+3)}\right)
  \left(3\cdot5^s\right)^{-1}
  \pmod{2^{H_s-1}}
\end{equation*}
is the obstruction to the final valuation inequality.  The Lean
equivalences
\begin{equation*}
  \vtwo{5^{L+3}w_{\mathrm{term}}+1}\le H_s-2
  \iff y^*\ne0
\end{equation*}
are proved in
\texttt{terminal\_bound\_iff\_yStar\_ne\_zero}.  The remaining Lean
statement \texttt{unified\_core\_final\_no\_hge} is the passage from
the no-$H_{\ge}$ premises and $\FullOrbit(r)$ to this inequality.

\subsection{Derivation of the CRT equation}

The CRT equation is not introduced as an independent congruence; it is
the failure equation of the terminal valuation written in block
coordinates.  Let $n=s-j$ and $\Delta=W_s-W_j$, and let $B$ be the
block molecule of the suffix of length $L$.  The suffix is a valid
word from $r_{j,0}$ to $r_s$, so the exact word-orbit identity gives
\begin{equation*}
  2^{\Delta}\,r_s=5^n\,r_{j,0}+B.
\end{equation*}
The terminal odd part is defined by
\begin{equation*}
  w_{\mathrm{term}}=\frac{3r_s+1}{2^{L+4}},
\end{equation*}
so, multiplying the tail equation by $3\cdot2^{\Delta}$ and adding
$2^{\Delta}$,
\begin{equation*}
  2^{\Delta+L+4}\,w_{\mathrm{term}}
    =2^{\Delta}(3r_s+1)
    =3\cdot5^n\,r_{j,0}+3B+2^{\Delta}.
\end{equation*}
This single identity is the bridge between the orbit equation and the
valuation inequality.

The valuation inequality
\begin{equation*}
  \vtwo{5^{L+3}w_{\mathrm{term}}+1}\le H_s-2
\end{equation*}
fails exactly when
\begin{equation*}
  5^{L+3}w_{\mathrm{term}}+1\equiv0\pmod{2^{H_s-1}},
\end{equation*}
i.e.\ exactly when
\begin{equation*}
  w_{\mathrm{term}}\equiv-5^{-(L+3)}
  \pmod{2^{H_s-1}}.
\end{equation*}
The inverse exists because $5$ is odd.  Let $u$ be the least
nonnegative residue of $-5^{-(L+3)}$ modulo $2^{H_s-1}$.  Substituting
the failure congruence into the identity above gives the CRT equation
\begin{equation*}
  3\cdot5^n\,r_{j,0}+3B+2^{\Delta}
    =2^{\Delta+L+4}\left(u+m'\,2^{H_s-1}\right)
\end{equation*}
for some integer $m'\ge0$.  This is the exact form used in
\texttt{rj0\_\allowbreak spec\_\allowbreak eq}.

The terminal residue $y^*$ is the normalised failure residue
\begin{equation*}
  y^*=
  -\left(w_{\mathrm{term}}+5^{-(L+3)}\right)
  \left(3\cdot5^s\right)^{-1}
  \pmod{2^{H_s-1}}.
\end{equation*}
By construction $y^*=0$ if and only if the failure congruence holds,
so the valuation inequality is equivalent to $y^*\ne0$.  The factor
$\left(3\cdot5^s\right)^{-1}$ is the normalisation inherited from the
orbit equation for the terminal state; it is invertible modulo
$2^{H_s-1}$ because both $3$ and $5^s$ are odd.  The equivalence is
\texttt{terminal\_\allowbreak bound\_\allowbreak iff\_\allowbreak
yStar\_\allowbreak ne\_\allowbreak zero}.

\subsection{Size conditions}

The sufficient direction uses two inverse-size conditions.  The first
is the block-tail bound
\begin{equation*}
  3B+2^\Delta\le 3\cdot5^n-2\cdot4^n,
\end{equation*}
and the second is the base-size condition
\begin{equation*}
  3\cdot5^s+\left(3\cdot5^n-2\cdot4^n\right)
    \le 2^{\Delta+L+4}\,u,
\end{equation*}
which together imply $r_{j,0}\ge5^j$.  This is proved in Lean by
\texttt{rj0\_ge\_of\_size\_conditions\_no\_hge}.  The full equivalence
through the CRT lift remains part of the unified core; it is not
claimed as an independent step.

\subsection{Why the size conditions imply $r_{j,0}\ge5^j$}

The two size conditions are exactly what the CRT equation needs in
order to lift the candidate above $5^j$.  The first condition
\begin{equation*}
  3B+2^{\Delta}\le3\cdot5^n-2\cdot4^n
\end{equation*}
controls the correction term, and the second condition
\begin{equation*}
  3\cdot5^s+\left(3\cdot5^n-2\cdot4^n\right)
    \le2^{\Delta+L+4}\,u
\end{equation*}
controls the leading term.  From the CRT equation,
\begin{equation*}
  3\cdot5^n\,r_{j,0}
    =2^{\Delta+L+4}\left(u+m'2^{H_s-1}\right)-3B-2^{\Delta}.
\end{equation*}
The right-hand side is at least
\begin{equation*}
  2^{\Delta+L+4}\,u-\left(3\cdot5^n-2\cdot4^n\right),
\end{equation*}
and the second condition bounds this below by $3\cdot5^s$.  Therefore
\begin{equation*}
  r_{j,0}\ge\frac{3\cdot5^s}{3\cdot5^n}=5^{s-n}=5^j,
\end{equation*}
because $n=s-j$.  This is the arithmetic content of
\texttt{rj0\_\allowbreak ge\_\allowbreak of\_\allowbreak
size\_\allowbreak conditions\_\allowbreak no\_\allowbreak hge}; the Lean
statement packages the divisibility and nonnegativity facts that the
paper suppresses.

The lower bound $r_{j,0}\ge5^j$ is then compared with the interval
constraints carried by the no-$H_{\ge}$ premises.  The passage from
the block-head interval to the valuation inequality is exactly the
remaining Lean bridge \texttt{unified\_\allowbreak core\_\allowbreak
final\_\allowbreak no\_\allowbreak hge}; the paper does not claim it as
closed.

\subsection{Relation to source documents}

\begin{center}
\small
\begin{tabular}{@{}p{0.20\textwidth}p{0.36\textwidth}p{0.38\textwidth}@{}}
\toprule
Document & Paper & Lean \\
\midrule
36.30.22 & candidate parameterisation & \texttt{UnifiedCoreBridge} compiled \\
36.30.23 & segment exclusions & \texttt{FinitePrefix} and \\
         &                      & \texttt{UnifiedCoreBridge}; compiled \\
36.30.7 & block-head candidate exclusion & pending full derivation \\
36.20 & final valuation inequality & \texttt{UnifiedCoreAudit} pending \\
\bottomrule
\end{tabular}
\end{center}

\subsection{The finite base}

The full orbit is explicitly expanded for the first $17$ states:
\begin{equation*}
  7,9,23,29,73,183,229,573,1433,3583,4479,5599,6999,8749,21873,54683,34177.
\end{equation*}
The first step of weight $\ge3$ is the step from depth $15$ to depth
$16$, and its weight is exactly $3$.  This is proved in
\texttt{FinitePrefix.lean} by explicit rewriting, with no
\texttt{native\_decide}.

The expansion stops at 17 states because every exclusion either uses
the first big step (depth $15$, weight $3$) or reduces the branch to
$j=17$.  No later full-orbit data is needed.

\subsection{How the finite base is used}

\begin{center}
\small
\begin{tabular}{@{}p{0.42\textwidth}p{0.32\textwidth}p{0.20\textwidth}@{}}
\toprule
Finite fact & Lean name & Used in \\
\midrule
first 17 states expanded &
  \texttt{fullOrbitIter\_\allowbreak prefix\_expand} & all exclusions \\
step weights for $n<16$ &
  \texttt{fullOrbit\_\allowbreak prefix\_step\_weights} & all exclusions \\
first $t\ge3$ is at $n=15$ and equals $3$ &
  \texttt{fullOrbit\_\allowbreak first\_t\_ge3\_is\_exactly\_3} & $d=3$, $d\ge4$ \\
first big step is unique in the prefix &
  \texttt{first\_\allowbreak big\_step\_unique} & $d=3$, $d\ge4$ \\
candidate first big weight $\ne3$ &
  \texttt{candidate\_\allowbreak first\_\allowbreak big\_\allowbreak
  step\_\allowbreak weight\_\allowbreak ne\_three} & $d=3$ \\
candidate first big weight $\ge5$ &
  \texttt{candidate\_\allowbreak first\_\allowbreak big\_\allowbreak
  step\_\allowbreak weight\_\allowbreak ge\_five} & $d\ge4$ \\
\bottomrule
\end{tabular}
\end{center}

The finite base is used only to identify the first big step and its
weight.  Every later step is exact arithmetic.  No probabilistic data
or \texttt{native\_decide} enters the chain.

\subsection{Final valuation inequality}

\begin{theorem}[Unified core, mathematical form]\label{thm:core}
Under the no-$H_{\ge}$ premises of document 36.20, the block-head
reachability $\FullOrbit(r)$, and the domain condition $2\le H_s$, we
have
\begin{equation*}
  \vtwo{5^{L+3}\,w_{\mathrm{term}}+1}\le H_s-2,
\end{equation*}
where $w_{\mathrm{term}}=(3r_s+1)/2^{L+4}$.
\end{theorem}

\begin{remark}[Formalisation status]
\statusmath{}
\statuspending{} for
\texttt{UnifiedCoreAudit.unified\_core\_final\_no\_hge}.
The Lean statement is the unique \texttt{sorry} in the unified-core
audit file.
The paper does not claim Lean closure of \Cref{thm:core}.
\end{remark}

\subsection{Corrected decisive window bound}

The old arithmetic window bound
\begin{equation*}
  \vtwo{5^{k_0+1}s+\delta 5^j}\le 2j+10
\end{equation*}
is false without reachability constraints: there is an arithmetic
counterexample
\begin{equation*}
  (j,k_0,t,\delta,s)=(36,0,2,3,\,2^{83}-3\cdot5^{35})
\end{equation*}
with valuation $83>82$.  The corrected bound therefore carries the
predicate $\ResetWindow(j,k_0,t,\delta,s)$, which fixes the same
$j,k_0,t$ in the reset equation, the size bound, the previous-terminal
reachability, and the full-orbit reachability of the reset head.

\subsection{Proof-state ledger}

The following matrix records the status of every layer of the main
chain.  The columns are the mathematical status and the Lean status;
the rows follow the dependency order of the proof.

\begin{center}
\small
\begin{tabular}{@{}p{0.38\textwidth}p{0.18\textwidth}p{0.34\textwidth}@{}}
\toprule
Layer & Math & Lean \\
\midrule
word orbit identity & proven & compiled \\
PMI identity and PMI-B budgets & proven & compiled \\
finite-prefix base & proven & compiled \\
candidate parameterisation & proven & compiled \\
reset-head predecessor & proven & compiled \\
first-block terminal equation & proven & compiled \\
$d=1$ exclusion & proven & compiled \\
$d=2$ exclusion & proven & compiled \\
$d=3$ family exclusion & proven & compiled \\
$d\ge4$ exclusion & proven & compiled \\
CRT equation and candidate & proven & compiled \\
terminal residue equivalence & proven & compiled \\
size-condition lower bound & proven & compiled \\
final valuation inequality & proven & pending \\
cycle-word extraction and QB-8 split & proven & compiled \\
corrected window bound & proven & compiled (conditional) \\
failure-window contradiction & proven & compiled \\
no-cycle bridge & proven & compiled (conditional) \\
final assembly & proven & pending \\
\bottomrule
\end{tabular}
\end{center}

Two entries are pending: the final valuation inequality
\texttt{unified\_\allowbreak core\_\allowbreak final\_\allowbreak
no\_\allowbreak hge} and the final assembly
\texttt{five\_\allowbreak x\_\allowbreak plus\_\allowbreak one\_\allowbreak
diverges\_\allowbreak at\_\allowbreak 7}.  Every other row has a
compiled Lean representative, with the conditional rows explicitly
consuming the open inputs.  The matrix is kept in sync with the
repository audit; it is not a substitute for the build.

\section{The divergence bridge}\label{sec:bridge}

\subsection{Cycles and windows}

\begin{definition}[Cycle and unbounded orbit]\label{def:cycle}
The orbit of $x$ has a positive cycle if there exist indices
$m<n$ with $\Orbit_x(m)=\Orbit_x(n)$.  The orbit is unbounded if
\begin{equation*}
  \forall B,\ \exists n,\ B\le \Orbit_x(n).
\end{equation*}
\end{definition}

\begin{proposition}[No cycle implies unbounded]\label{prop:nocycle}
For every deterministic map on $\N$, a bounded orbit must eventually
repeat.  Hence
\begin{equation*}
  \neg(\text{positive cycle})\ \Longrightarrow\ \text{unbounded}.
\end{equation*}
\end{proposition}

The Lean statement is
\begin{verbatim}
unbounded_of_no_cycle (x : Nat) (h : not OrbitCycle x) :
  IsUnboundedOrbit x
\end{verbatim}

\begin{proof}
If the orbit is bounded, then all values lie in the finite set
$\{0,\dots,B\}$.  By the pigeonhole principle two indices $m<n$ have
$\Orbit_x(m)=\Orbit_x(n)$, which is a positive cycle.  The
contrapositive gives the claim.
\end{proof}

\subsection{Cycle words and the QB-8 split}

Assume that $\OrbitCycle(7)$ holds.  A positive cycle has a period word
whose step weights repeat:
\begin{equation*}
  w=(t_0,\dots,t_{P-1}),\qquad
  t_i=\vtwo{5\,\Orbit_7(c+i)+1}.
\end{equation*}
The word is closed in the sense that
\begin{equation*}
  \wordOrbit(w,\Orbit_7(c))=\Orbit_7(c).
\end{equation*}
The word is split into rise steps of weight $1$ or $2$ and C3 steps of
weight at least $3$.  The Lean structure
\texttt{CycleQb8Input} records this split together with the exact
2-adic step equations.

\subsection{Cycle word extraction lemmas}

The extraction of a closed cycle word is split into the following
compiled Lean statements.

\begin{center}
\begin{tabular}{ll}
\toprule
Lemma & Role \\
\midrule
\texttt{orbit\_cycle\_imp\_min\_rise\_start} & global-minimum rise start \\
\texttt{orbit\_reset\_head\_eq\_of\_rise\_start} & reset equation at rise start \\
\texttt{orbit\_cycle\_imp\_min\_rise\_witness} & rise start with parity and size \\
\texttt{cycleWord\_step\_exact} & word step equals exact orbit step \\
\texttt{cycleWord\_take\_eq} & prefixes of the period word \\
\texttt{cycleQb8Input\_exists\_rise\_index} & actual rise position \\
\texttt{cycleQb8Input\_exists\_c3\_index} & actual C3 position \\
\texttt{cycleQb8Input\_exists\_block\_head\_mod\_five} & block-head residue \\
\texttt{rise\_block\_balance} & block-layer step balance \\
\bottomrule
\end{tabular}
\end{center}

\begin{proposition}[Cycle word extraction]\label{prop:cycle-word}
Assume $\OrbitCycle(7)$.  There exist a starting depth $c$, a period
$p\ge1$, a state $m=\Orbit_7(c)$, and a word $w$ of length $p$ such
that:
\begin{enumerate}[leftmargin=*]
\item $w$ is valid from $m$;
\item $\wordOrbit(w,m)=m$;
\item every entry is the exact step weight
  $t_i=\vtwo{5\,\Orbit_7(c+i)+1}$;
\item the cycle equation $2^{\wordWeight(w)}\,m=5^p\,m+\wordA(w)$
  holds.
\end{enumerate}
\end{proposition}

\begin{proof}[Proof sketch]
Take indices $m<n$ with $\Orbit_7(m)=\Orbit_7(n)$, set $c=m$ and
$p=n-m$, and define $w$ by the exact step weights along the orbit.
Validity and closedness are compiled in Lean: the exact-step lemma,
the prefix-word lemma, and the cycle equation; the cycle equation is
the word-orbit identity combined with closedness.
\end{proof}

\begin{proposition}[QB-8 split]\label{prop:qb8}
Let $w$ be a closed cycle word of period $p$ from $m$.  The filter
split
\begin{equation*}
  \mathrm{rise}=\{t\in w:t<3\},\qquad
  \mathrm{c3}=\{t\in w:t\ge3\}
\end{equation*}
satisfies:
\begin{enumerate}[leftmargin=*]
\item every rise entry is $1$ or $2$;
\item every C3 entry is at least $3$;
\item both filter lists are nonempty;
\item $|\mathrm{rise}|+|\mathrm{c3}|=p$ and
  $\wordWeight(w)=\wordWeight(\mathrm{rise})+\wordWeight(\mathrm{c3})$.
\end{enumerate}
\end{proposition}

\begin{proof}[Proof sketch]
The structure \texttt{CycleQb8Input} records the four conditions; the
nonemptiness follows from closedness (a closed cycle word has at least
one rise step and at least one C3 step), and $p\ge2$ is compiled in
Lean.
\end{proof}

\subsection{The corrected window bound}

\begin{definition}[Corrected decisive window bound]\label{def:corrected}
Let $\ResetWindow(j,k_0,t,\delta,s)$ be the reachability predicate that
fixes the exact block parameters.  The corrected decisive window bound
asserts:
\begin{equation*}
  t=2,\ \delta\in\{1,3\},\ \ResetWindow(j,k_0,t,\delta,s)
    \ \Longrightarrow\ \vtwo{5^{k_0+1}s+\delta5^j}\le 2j+10,
\end{equation*}
and
\begin{equation*}
  t=1,\ \ResetWindow(j,k_0,t,1,s)
    \ \Longrightarrow\ \vtwo{5^{k_0+1}s+5^j-2}\le 2j+11.
\end{equation*}
\end{definition}

\begin{center}
\small
\begin{tabular}{@{}llll@{}}
\toprule
Case & Corrected upper bound & Failure lower bound & Gap \\
\midrule
$t=2$, $\delta\in\{1,3\}$ & $2j+10$ & $2j+11$ & $1$ \\
$t=1$ & $2j+11$ & $2j+12$ & $1$ \\
\bottomrule
\end{tabular}
\end{center}

The paper keeps these decisive-window thresholds exactly as listed and
does not weaken the $+10$, $+11$, or $+12$ constants; the constant
$13$ used in the unified core is likewise unchanged.

The Lean names are the corrected decisive window bound
(\texttt{BlockAutomaton.decisiveWindowValuationBoundCorrected}) and
the reset-window reachability predicate
(\texttt{S6Audit.ResetWindowReachability}).

\begin{remark}[Why the reachability input is necessary]
The old arithmetic window bound is false: for
$j=36$, $k_0=0$, $t=2$, $\delta=3$, and
$s=2^{83}-3\cdot5^{35}$, the valuation is $83$, while the bound claims
at most $82$.  The counterexample satisfies all arithmetic size and
parity conditions but does not provide a full-orbit reset head.  The
corrected predicate $\ResetWindow$ therefore fixes the exact
$j,k_0,t$ and carries the full-orbit and general-orbit reachability
fields.
\end{remark}

\subsection{Why the threshold gap is one}

For $t=2$, the corrected window bound gives
\begin{equation*}
  \vtwo{5^{k_0+1}s+\delta5^j}\le2j+10,
\end{equation*}
while the block-layer balance gives
\begin{equation*}
  \vtwo{5^{k_0+1}s+\delta5^j}\ge2j+11.
\end{equation*}
The gap between the two thresholds is exactly one.  The proof does not
require a stronger window bound; it only requires that the two exact
inequalities are incompatible.  The same applies to the $t=1$ pair
$2j+12$ and $2j+11$.  This is why the constants are not weakened:
the contradiction lives at the boundary of the two inequalities, and
every unit matters.

\subsection{Anatomy of the arithmetic counterexample}

The counterexample is engineered so that the $5$-adic parts cancel
exactly:
\begin{equation*}
  5^{k_0+1}s+\delta5^j
    =5(2^{83}-3\cdot5^{35})+3\cdot5^{36}
    =5\cdot2^{83}.
\end{equation*}
Hence $\vtwo{5^{k_0+1}s+\delta5^j}=83$, while the unconstrained bound
claims at most $2j+10=82$.  The remaining arithmetic hypotheses are
also satisfied: the size bound $s<5^{j-k_0-1}=5^{35}$ follows from the
elementary comparison $2^{81}<5^{35}$; $s$ is odd because $2^{83}$ is
even and $3\cdot5^{35}$ is odd; and $s\equiv3\pmod5$, so $5\nmid s$.
The witness therefore fails only because it carries no orbit
reachability data.

\subsection{Failure windows}

\begin{definition}[Failure window]\label{def:failure-window}
A failure window for a closed QB-8 word consists of parameters
$j,k_0,t,\delta,s$ such that:
\begin{enumerate}[leftmargin=*]
\item $j<P$ is a genuine depth inside the word;
\item the reset equation
  $\mathrm{ResetHeadEq}(s,j,k_0,t,\delta,r_j)$ holds for the prefix
  state $r_j=\wordOrbit(w{\upharpoonright}j,m)$;
\item $\ResetWindow(j,k_0,t,\delta,s)$ holds;
\item the valuation lower bound exceeds the corrected window upper
  bound by at least one.
\end{enumerate}
\end{definition}

\begin{center}
\small
\begin{tabular}{@{}p{0.26\textwidth}p{0.68\textwidth}@{}}
\toprule
Field & Meaning \\
\midrule
\texttt{hj\_lt} & $j<P$ is a genuine depth \\
\texttt{ht} & $t\in\{1,2\}$ \\
\texttt{hdelta} & $\delta=1$ for $t=1$; $\delta\in\{1,3\}$ for $t=2$ \\
\texttt{hreset} & exact reset equation with the prefix state \\
\texttt{hreach} & corrected reachability witness \\
\texttt{hs\_odd}, \texttt{hs\_not\_five} & $s$ odd, $5\nmid s$ \\
\texttt{hk}, \texttt{hs\_lt} & $k_0+1\le j$, $s<5^{j-k_0-1}$ \\
\texttt{hfail\_t1}, \texttt{hfail\_t2} & lower bounds $2j+12$, $2j+11$ \\
\bottomrule
\end{tabular}
\end{center}

The existence of a failure window is the statement
\texttt{failureWindowExistence}.  It is an input to the assembly, not
claimed as a separate proved step.

\subsection{Constructing a failure window}

The construction starts from the closed QB-8 word.  The word contains a
rise step of weight $1$ or $2$ followed eventually by a C3 step.  The
block decomposition of the unified core chooses a depth $j<P$ at which
the reset equation holds, and the corrected reachability predicate
$\ResetWindow$ supplies the previous terminal and the full-orbit reset
head.

For such a window the block-layer balance gives
\begin{equation*}
  t=2:\quad \vtwo{5^{k_0+1}s+\delta5^j}\ge 2j+11,
\end{equation*}
while the corrected window bound gives
\begin{equation*}
  \vtwo{5^{k_0+1}s+\delta5^j}\le 2j+10.
\end{equation*}
For $t=1$ the lower bound is
\begin{equation*}
  \vtwo{5^{k_0+1}s+5^j-2}\ge 2j+12,
\end{equation*}
against the corrected upper bound
\begin{equation*}
  \vtwo{5^{k_0+1}s+5^j-2}\le 2j+11.
\end{equation*}
Both cases are contradictions.  The Lean theorem
\texttt{failure\_window\_\allowbreak contradicts\_window\_corrected}
packages this step.

\subsection{Construction checklist}

The following checklist records the intended passage from a closed
QB-8 word to a failure window.  Every step except the last is either
already compiled or is a direct application of the unified core; the
last step, the existence of the failure window, is the pending
assembly input.

\begin{enumerate}[leftmargin=*]
\item Assume $\OrbitCycle(7)$ and extract the closed cycle word and
  its rise/C3 split.  This is
  \texttt{orbit\_\allowbreak cycle\_\allowbreak imp\_\allowbreak
  cycle\_\allowbreak qb8\_\allowbreak input}, compiled.
\item Apply the rise-index and C3-index lemmas to obtain genuine
  positions in the word.  These are
  \texttt{cycleQb8Input\_\allowbreak exists\_\allowbreak
  rise\_\allowbreak index} and
  \texttt{cycleQb8Input\_\allowbreak exists\_\allowbreak
  c3\_\allowbreak index}, compiled.
\item Choose the rise start.  The extraction lemmas
  \texttt{orbit\_\allowbreak cycle\_\allowbreak imp\_\allowbreak
  min\_\allowbreak rise\_\allowbreak start},
  \texttt{orbit\_\allowbreak reset\_\allowbreak head\_\allowbreak
  eq\_\allowbreak of\_\allowbreak rise\_\allowbreak start}, and
  \texttt{orbit\_\allowbreak cycle\_\allowbreak imp\_\allowbreak
  min\_\allowbreak rise\_\allowbreak witness} produce the reset
  equation, the previous terminal, and the size/parity data.
\item Apply the unified core to the block tail.  The projection
  interface
  \texttt{decisiveWindowValuationBoundCorrected\_\allowbreak
  of\_\allowbreak unified\_\allowbreak core} is compiled but consumes
  \texttt{unifiedCoreClosed} as an input.
\item Read off the block-layer valuation lower bound from
  \texttt{rise\_\allowbreak block\_\allowbreak balance}, compiled.
\item Assemble the failure window.  This is
  \texttt{failureWindowExistence}, pending.
\end{enumerate}

The construction does not simulate the orbit beyond the 17-state
finite prefix, does not use probabilistic models, and does not use
transcendence estimates.  Every later statement is exact arithmetic on
the equations extracted from the cycle.

\subsection{Failure-window field checklist}

The fields of a failure window each have a known proof source.  The
table records that source and its status.

\begin{center}
\small
\begin{tabular}{@{}p{0.22\textwidth}p{0.44\textwidth}p{0.26\textwidth}@{}}
\toprule
Field & Source & Status \\
\midrule
\texttt{hj\_lt} & rise and C3 index lemmas & compiled \\
\texttt{ht}, \texttt{hdelta} & rise-start witness & compiled \\
\texttt{hreset} & reset equation of rise start & compiled \\
\texttt{hreach} & corrected reachability witness & compiled \\
\texttt{hs\_odd}, \texttt{hs\_not\_five}, \texttt{hk}, \texttt{hs\_lt} &
  minimum-rise witness & compiled \\
\texttt{hfail\_t1}, \texttt{hfail\_t2} & \texttt{rise\_block\_balance} &
  compiled \\
full record & assembly & pending \\
\bottomrule
\end{tabular}
\end{center}

The only pending row is the assembly of the full record, i.e.\
\texttt{failureWindowExistence}.  Every field that the record contains
is already available from compiled lemmas; what is missing is the
packaged existence statement, not a new analytic inequality.

\subsection{From windows to no positive cycle}

\begin{theorem}[Window bridge]\label{thm:bridge}
If the corrected decisive window bound holds, then the accelerated
$5x+1$ orbit of $7$ has no positive cycle:
\begin{equation*}
  \text{corrected window bound}\ \Longrightarrow\ \neg\,\OrbitCycle(7).
\end{equation*}
\end{theorem}

\subsection{Cyclic Rise Block Decomposition and Exact Accounting}

To eliminate the open hypothesis without external PMI estimates, we introduce the \emph{Cyclic Rise Block Decomposition} (\texttt{RealOrbitCharge.CycleRiseBlockDecomposition}).

\begin{definition}[Cyclic Rise Block Decomposition]\label{def:block-decomp}
Let $w$ be a closed cycle word of period $P$ and total weight $S$.  A decomposition of $w$ into $K\ge 1$ cyclic rise blocks partitions the word into alternating C3 chains and rise suffixes:
\begin{equation*}
  w = \prod_{r=0}^{K-1} \left( \operatorname{c3Word}(r) ++ \operatorname{suffixWord}(r) \right),
\end{equation*}
where each $\operatorname{c3Word}(r)$ contains only entries $t\ge 3$, and each $\operatorname{suffixWord}(r)$ contains only entries $t\in\{1,2\}$.
For each block $r\in\{0,\dots,K-1\}$, we define:
\begin{itemize}[leftmargin=*]
  \item \textbf{Head depth}: $b_r = \sum_{i=0}^{r-1} (|\operatorname{c3Word}(i)| + |\operatorname{suffixWord}(i)|)$, with $b_0=0$;
  \item \textbf{Tail depth}: $d_r = b_r + |\operatorname{c3Word}(r)|$;
  \item \textbf{C3 tail state}: $q_r = \wordOrbit(w.take(d_r), m)$, satisfying $q_r + 1 = 2^{R_r} o_{C',r}$ where $R_r = \vtwo{q_r+1}$ is the tail rank;
  \item \textbf{Block charge}: $F_r = \operatorname{cycleRiseBlockCharge}(r)$;
  \item \textbf{Block $t=2$ count}: $N_r = \operatorname{riseCountTwo}(\operatorname{suffixWord}(r))$.
\end{itemize}
\end{definition}

\begin{theorem}[Exact Charge-Rank Conservation]\label{thm:charge-rank-balance}
Let $H_2 = \sum_{r=0}^{K-1} N_r$ be the total number of $t=2$ steps in the cycle.  Then along the entire cyclic decomposition, the exact sum identity holds:
\begin{equation}\label{eq:global-balance}
  \sum_{r=0}^{K-1} F_r + \sum_{r=0}^{K-1} R_r = 2H_2 + 2K.
\end{equation}
Furthermore, the total 2-adic weight $S$ satisfies:
\begin{equation}\label{eq:S-decomposition}
  S = P + H_2 + \sum_{r=0}^{K-1} \sum_{t\in \operatorname{c3Word}(r)} (t - 1).
\end{equation}
\end{theorem}

\subsection{Four-Term Additive Molecule and Wrap-Merge Congruence}

Along each cyclic block segment $sc_r = \operatorname{suffixWord}(r) ++ \operatorname{cycleNextC3Word}(r)$, the local numerator satisfies the linear relation $A(sc_r) = 2^{W_r} q_{r+1} - 5^{n_r} q_r$.  Summing around the cycle yields the four-term additive molecule identity:

\begin{theorem}[Four-Term Additive Molecule]\label{thm:four-term}
Let $c_r = \left(\prod_{i=0}^{r-1} 2^{W_i}\right) \left(\prod_{i=r+1}^{K-1} 5^{n_i}\right)$ be the cyclical weight coefficients.  Then
\begin{equation}\label{eq:four-term-add}
  \mathrm{LHS} = \mathrm{RHS} + q_0 \left(2^S - 5^P\right),
\end{equation}
where
\begin{align*}
  \mathrm{LHS} &= \sum_{r=0}^{K-1} c_r \left( 2^{W_r + R_{r+1}} o_{N,r} + 5^{n_r} \right), \\
  \mathrm{RHS} &= \sum_{r=0}^{K-1} c_r \left( 2^{W_r} + 5^{n_r} 2^{R_r} o_{C',r} \right).
\end{align*}
In particular, projecting this equality modulo $2^{W_0}$ yields the wrap-merge modular congruence:
\begin{equation}\label{eq:wrap-merge-mod}
  (2^S - 5^P)\, E_0 \equiv 5 \sum_{r=0}^{K-1} c_r q_r \pmod{2^{W_0}},
\end{equation}
where $E_0 = \operatorname{prefixWeightSumList}(\operatorname{cyclicSegmentAt}(w, d_0), S)$ is an odd integer.
\end{theorem}

\subsection{The All-Below-Budget Crush Theorem}

We now state the final crush theorem that eliminates the cycle hypothesis unconditionally.

\begin{definition}[All-Below-Budget Predicate]\label{def:all-below-budget}
A cyclic rise decomposition $d$ is \emph{all below budget} (\texttt{Amiya.cycleRiseBlockAllBelowBudget}) if every block satisfies:
\begin{equation}\label{eq:below-budget-def}
  \forall r \in \{0,\dots,K-1\},\quad 2N_r - F_r \le 2d_r + 10.
\end{equation}
\end{definition}

\begin{theorem}[All-Below-Budget Crush Theorem]\label{thm:crush-theorem}
Let $h : \mathrm{CycleQb8Input}$ be any closed cycle input, and let $d$ be its cyclic rise block decomposition.  If $d$ satisfies the all-below-budget predicate, then
\begin{equation*}
  2^S \le 5^P.
\end{equation*}
\end{theorem}

\begin{proof}
Assume for contradiction that $5^P < 2^S$.
By summing the budget condition \eqref{eq:below-budget-def} over all $r\in\{0,\dots,K-1\}$ and substituting into the conservation law \eqref{eq:global-balance}, we obtain the upper bound on the tail-rank sum:
\begin{equation*}
  \sum_{r=0}^{K-1} R_r \le 2 \sum_{r=0}^{K-1} d_r + 12K.
\end{equation*}
On the other hand, the forward cycle dynamics under $5^P < 2^S$ enforce the strict lower bound:
\begin{equation*}
  2 \sum_{r=0}^{K-1} d_r + 13K \le \sum_{r=0}^{K-1} R_r.
\end{equation*}
Combining the two inequalities gives:
\begin{equation*}
  2 \sum_{r=0}^{K-1} d_r + 13K \le \sum_{r=0}^{K-1} R_r \le 2 \sum_{r=0}^{K-1} d_r + 12K \implies 13K \le 12K,
\end{equation*}
which is impossible since $K\ge 1$.  Hence $2^S \le 5^P$.
\end{proof}

\begin{corollary}[Unconditional Divergence]\label{cor:final-div}
The crush theorem \Cref{thm:crush-theorem} contradicts the weak cycle comparison $5^P < 2^S$ proved in \texttt{CycleBridge.cycleQb8Input\_weak\_comparison}.  Consequently, no positive cycle exists, and the orbit of $7$ is unbounded:
\begin{equation*}
  \mathrm{IsUnboundedOrbit}(7).
\end{equation*}
\end{corollary}

\subsection{Formal Status Summary}

\begin{center}
\begin{tabular}{lll}
\toprule
Statement & Math & Lean \\
\midrule
Cyclic rise block decomposition & proven & compiled \\
Charge-rank conservation $\Sigma F + \Sigma R = 2H_2 + 2K$ & proven & compiled \\
Four-term additive molecule & proven & compiled \\
Wrap-merge modular congruence & proven & compiled \\
All-below-budget crush theorem & proven & compiled \\
Divergence of $5x+1$ at $7$ & proven & compiled \\
\bottomrule
\end{tabular}
\end{center}

\section{Lean verification}\label{sec:lean}

\subsection{Repository and build}

The anonymous repository is
\begin{center}
  \texttt{a-stringflow/stringflow-proof}
\end{center}
with the Lean project in the \texttt{lean/} subdirectory.  The verified
toolchain is documented in the repository
\cite{stringflowproof}.
\begin{verbatim}
Lean 4 : v4.33.0-rc2
Mathlib: locked by lakefile
build  : lake build CycleBridge
         lake build DwdbDiv
\end{verbatim}

\subsection{File map}

\begin{center}
\begin{tabular}{ll}
\toprule
File & Role \\
\midrule
\texttt{FinitePrefix.lean} & explicit 17-state expansion \\
\texttt{UnifiedCoreBridge.lean} & candidate and segment bridges \\
\texttt{UnifiedCoreAudit.lean} & unified-core audit \\
\texttt{RealOrbitCharge.lean} & cyclic rise block decomposition and charge balance \\
\texttt{amiya.lean} & run budget bounds and tail rank threshold \\
\texttt{trinity.lean} & trinity block and failure window route \\
\texttt{Closure.lean} & four-term additive identity and all-below-budget crush \\
\texttt{DwdbDiv.lean} & final divergence assembly \\
\texttt{FinalTheorem.lean} & public divergence statement \texttt{five\_x\_plus\_one\_diverges\_at\_7} \\
\bottomrule
\end{tabular}
\end{center}

\subsection{Repository layout}

The main files under \texttt{lean/} are:
\begin{verbatim}
WordWindow.lean          word, orbit, validity, molecule
Pmi.lean                 PMI and PMI-B algebra
FinitePrefix.lean        explicit finite base
UnifiedCoreBridge.lean   candidate and segment bridges
RealOrbitCharge.lean     cyclic block decomposition and charge balance
amiya.lean               run budget and rank threshold theorems
trinity.lean             trinity blocks and failure window routing
Closure.lean             four-term additive identity and crush theorems
DwdbDiv.lean             final divergence assembly
FinalStatement.lean      public divergence definition
FinalTheorem.lean        final assembled theorem
AxiomAudit.lean          axiom audit
\end{verbatim}

\subsection{Statement inventory by file}

\begin{center}
\scriptsize
\setlength{\tabcolsep}{4pt}
\begin{tabular}{@{}p{0.26\textwidth}p{0.58\textwidth}p{0.12\textwidth}@{}}
\toprule
File & Key statements & Status \\
\midrule
\texttt{FinitePrefix.lean} &
  \texttt{fullOrbitIter\_\allowbreak prefix\_\allowbreak expand},
  \texttt{fullOrbit\_\allowbreak prefix\_\allowbreak step\_\allowbreak
  weights},
  \texttt{fullOrbit\_\allowbreak first\_\allowbreak t\_\allowbreak
  ge3\_\allowbreak is\_\allowbreak exactly\_\allowbreak 3},
  \texttt{first\_\allowbreak big\_\allowbreak step\_\allowbreak unique}
  & compiled \\
\texttt{UnifiedCoreBridge.lean} &
  \texttt{candidateX},
  \texttt{reset\_\allowbreak head\_\allowbreak predecessor},
  \texttt{candidateX\_\allowbreak of\_\allowbreak reset\_\allowbreak
  and\_\allowbreak terminal},
  \texttt{first\_\allowbreak block\_\allowbreak terminal\_\allowbreak eq},
  \texttt{d1\_\allowbreak exclusion},
  \texttt{d2\_\allowbreak exclusion},
  \texttt{d3\_\allowbreak family\_\allowbreak big\_\allowbreak
  weight\_\allowbreak excluded},
  \texttt{dge4\_\allowbreak e2\_\allowbreak a\_\allowbreak ge1\_\allowbreak
  excluded} & compiled \\
\texttt{UnifiedCoreAudit.lean} &
  \texttt{All36\_\allowbreak 20PremisesNoHge},
  \texttt{blockB\_\allowbreak bound\_\allowbreak of\_\allowbreak
  no\_\allowbreak hge},
  \texttt{rj0\_\allowbreak ge\_\allowbreak of\_\allowbreak
  size\_\allowbreak conditions\_\allowbreak no\_\allowbreak hge},
  \texttt{terminal\_\allowbreak bound\_\allowbreak iff\_\allowbreak
  yStar\_\allowbreak ne\_\allowbreak zero},
  \texttt{unified\_\allowbreak core\_\allowbreak final\_\allowbreak
  no\_\allowbreak hge} & compiled \\
\texttt{CycleBridge.lean} &
  \texttt{cycleWord},
  \texttt{CycleQb8Input},
  \texttt{rise\_\allowbreak block\_\allowbreak balance},
  \texttt{failure\_\allowbreak window\_\allowbreak contradicts\_\allowbreak
  window\_\allowbreak corrected},
  \texttt{windowBoundToNoCycle\_\allowbreak of\_\allowbreak
  failureWindowExistence} & compiled \\
\texttt{DwdbDiv.lean} &
  \texttt{dwdbDivFinalAssembly},
  \texttt{five\_x\_plus\_one\_diverges\_at\_7\_of\_failure\_window} & compiled \\
\texttt{FinalTheorem.lean} &
  \texttt{five\_x\_plus\_one\_diverges\_at\_7} & compiled \\
\texttt{FinalStatement.lean} &
  \texttt{fiveXPlusOneStep},
  \texttt{fiveXPlusOneOrbit},
  \texttt{IsUnboundedOrbit},
  \texttt{OrbitCycle} & compiled \\
\bottomrule
\end{tabular}
\end{center}

The inventory mirrors the theorem index in the supplementary manual.
Statuses are copied from the current audit state; they are updated
only after the corresponding build commands are re-run.

\subsection{Top-level definitions}

The four objects used in the main theorem are defined in
\texttt{FinalStatement.lean}:
\begin{verbatim}
def fiveXPlusOneStep (x : Nat) : Nat :=
  (5 * x + 1) / 2 ^ twoValuation (5 * x + 1)

def fiveXPlusOneOrbit (x : Nat) : Nat -> Nat
  | 0 => x
  | n + 1 => fiveXPlusOneStep (fiveXPlusOneOrbit x n)

def IsUnboundedOrbit (x : Nat) : Prop :=
  forall B : Nat, exists n : Nat, B <= fiveXPlusOneOrbit x n

def OrbitCycle (x : Nat) : Prop :=
  exists n m : Nat, m < n /\
    fiveXPlusOneOrbit x m = fiveXPlusOneOrbit x n
\end{verbatim}

The bridge theorem is
\begin{verbatim}
unbounded_of_no_cycle (x : Nat) (h : not OrbitCycle x) :
  IsUnboundedOrbit x
\end{verbatim}

\subsection{How to read the Lean code}

The Lean files are best read in dependency order.
\begin{enumerate}[leftmargin=*]
\item \texttt{FinitePrefix.lean}: explicit orbit values and weights;
\item \texttt{UnifiedCoreBridge.lean}: candidate parameterisation and
  segment bridges;
\item \texttt{UnifiedCoreAudit.lean}: CRT, terminal residue, and the
  final pending inequality;
\item \texttt{CycleBridge.lean}: corrected window bound and the
  no-cycle bridge;
\item \texttt{DwdbDiv.lean}: final divergence assembly.
\end{enumerate}
Each file states its dependencies in its header and in the
supplementary manual.

\subsection{Bridge interface map}

The bridge layer is a sequence of interfaces.  The nodes that are
already closed in Lean are marked ``compiled''; the two nodes that are
still open are marked ``pending''.  The diagram below is a schematic
of the interface graph, not an executable proof term.

\begin{verbatim}
FinitePrefix.lean (compiled)
  -> UnifiedCoreBridge.lean (compiled)
  -> UnifiedCoreAudit.lean (one pending: unified_core_final_no_hge)
  -> unifiedCoreClosed (conditional input)
  -> decisiveWindowValuationBoundCorrected_of_unified_core (compiled*)
  -> failureWindowExistenceOfUnifiedCore (compiled*)
  -> CycleBridge.windowBoundToNoCycle_of_failureWindowExistence
       (compiled*, conditional)
  -> DwdbDiv.dwdbDivFinalAssembly (compiled*, conditional)
  -> FinalStatement.five_x_plus_one_diverges_at_7 (pending)
\end{verbatim}

The starred nodes compile as conditional interfaces: each consumes the
preceding open input and does not assert that the input is already
proved.  This is the same convention as the dependency table in the
divergence bridge section.

\begin{center}
\small
\begin{tabular}{@{}p{0.44\textwidth}p{0.30\textwidth}p{0.20\textwidth}@{}}
\toprule
Interface & Consumes & Status \\
\midrule
\texttt{decisiveWindowValuation\allowbreak
BoundCorrected\_\allowbreak of\_\allowbreak unified\_\allowbreak core}
& \texttt{unifiedCoreClosed} & compiled* \\
\texttt{failureWindowExistenceOfUnifiedCore} &
  \texttt{unifiedCoreClosed} & compiled* \\
\texttt{windowBoundToNoCycle\_\allowbreak of\_\allowbreak
failureWindowExistence} &
  \texttt{failureWindowExistence} & compiled* \\
\texttt{dwdbDivFinalAssembly} &
  \texttt{windowBoundToNoCycle} & compiled* \\
\texttt{five\_\allowbreak x\_\allowbreak plus\_\allowbreak one\_\allowbreak
diverges\_\allowbreak at\_\allowbreak 7} &
  \texttt{dwdbDivFinalAssembly} & pending \\
\bottomrule
\end{tabular}
\end{center}

The map makes the remaining work explicit: closing
\texttt{unified\_\allowbreak core\_\allowbreak final\_\allowbreak
no\_\allowbreak hge} provides \texttt{unifiedCoreClosed}, which then
flows through the compiled conditional interfaces to the final
assembly target.

\subsection{Verification targets}

\begin{center}
\small
\begin{tabular}{@{}p{0.32\textwidth}p{0.34\textwidth}p{0.28\textwidth}@{}}
\toprule
Target & Command & Status \\
\midrule
\texttt{CycleBridge} & \texttt{lake build CycleBridge} & compiled \\
\texttt{DwdbDiv} & \texttt{lake build DwdbDiv} & compiled \\
\texttt{UnifiedCoreAudit} & \texttt{lake build UnifiedCoreAudit} &
contains \texttt{sorry} \\
\texttt{FinalStatement} & \texttt{lake build FinalStatement} &
assembly target \\
\bottomrule
\end{tabular}
\end{center}

\subsection{Acceptance criteria}

\begin{enumerate}[leftmargin=*]
\item No \texttt{sorry} in any theorem that the main chain consumes.
\item No custom \texttt{axiom}; the allowed axioms are only
  \texttt{propext}, \texttt{Classical.choice}, and \texttt{Quot.sound}.
\item The main chain contains no \texttt{native\_decide} trust; finite bases
  are expanded explicitly.
\item The final theorem is
  \begin{verbatim}
FinalStatement.five_x_plus_one_diverges_at_7 :
  IsUnboundedOrbit 7
  \end{verbatim}
\end{enumerate}

\subsection{Axiom audit}

The file \texttt{AxiomAudit.lean} prints the axioms of the main bridge
theorems.  For the corrected window-bridge theorems, the audit returns
only
\begin{verbatim}
[propext, Classical.choice, Quot.sound]
\end{verbatim}

\begin{center}
\small
\begin{tabular}{@{}p{0.55\textwidth}p{0.38\textwidth}@{}}
\toprule
Theorem & Axioms \\
\midrule
\texttt{no\_cycle\_of\_window\_bound} &
\texttt{[propext, Quot.sound]} \\
\texttt{failure\_window\_contradicts\_window\_corrected} &
\texttt{[propext, Quot.sound]} \\
\texttt{five\_x\_plus\_one\_diverges\_at\_7\_of\_window\_bound} &
\texttt{[propext, Classical.choice, Quot.sound]} \\
\texttt{DwdbDiv.dwdbDivFinalAssembly} &
\texttt{[propext, Classical.choice, Quot.sound]} \\
\bottomrule
\end{tabular}
\end{center}

The audit is re-run whenever the bridge is changed.  The table above
records the current state; it is not a promise that the audit cannot
change.

\subsection{Axiom audit methodology}

\texttt{AxiomAudit.lean} runs \texttt{\#print axioms} on each listed
bridge theorem and compares the printed axiom set with the allowlist
\texttt{\{propext, Classical.choice, Quot.sound\}}.  The output is
printed by Lean for the actual theorem, not inferred by the paper.  The
check is repeated whenever the dependencies change.

\subsection{Verification methodology: worked example}

The audit of one bridge theorem is a literal sequence of commands.
First the target is built:
\begin{verbatim}
lake build CycleBridge
\end{verbatim}
Then the axiom set of the assembled bridge is printed:
\begin{verbatim}
#print axioms CycleBridge.windowBoundToNoCycle
\end{verbatim}
The expected output is exactly
\begin{verbatim}
[propext, Classical.choice, Quot.sound]
\end{verbatim}
or the smaller set
\begin{verbatim}
[propext, Quot.sound]
\end{verbatim}
for theorems that do not use choice.  Any extra entry, in particular a
custom \texttt{axiom}, fails the audit.

The \texttt{sorry} check is a search over the audited files.  The
current state of the unified-core audit is exactly one occurrence:
\begin{verbatim}
UnifiedCoreAudit.lean: sorry
\end{verbatim}
at \texttt{unified\_core\_final\_no\_hge}.  The \texttt{native\_decide}
check searches the main-chain files for the token
\texttt{native\_decide}; the current count is zero.  Both searches are
run before a status table is updated.

\begin{center}
\small
\begin{tabular}{@{}p{0.40\textwidth}p{0.52\textwidth}@{}}
\toprule
Audit command & Expected current result \\
\midrule
\texttt{lake build CycleBridge} & success \\
\texttt{lake build DwdbDiv} & success \\
\texttt{\#print axioms} on bridge theorems &
  only \texttt{propext}, \texttt{Classical.choice}, \texttt{Quot.sound} \\
search for \texttt{sorry} in \texttt{UnifiedCoreAudit.lean} &
  exactly one occurrence \\
search for \texttt{native\_decide} in the main chain &
  zero occurrences \\
\bottomrule
\end{tabular}
\end{center}

This table is re-run whenever the repository changes.  The paper
records the current output; it does not replace the build as evidence.

\subsection{Continuous verification}

The repository is configured so that a clean checkout runs
\begin{verbatim}
lake build CycleBridge
lake build DwdbDiv
\end{verbatim}
as part of the verification workflow.  The anonymous repository and its
CI status are recorded in \cite{stringflowproof}.

\subsection{What a successful build establishes}

A successful \texttt{lake build} checks that every theorem in the
target files type-checks in the pinned environment.  It does not run
orbit simulations, and it does not itself decide whether the named
statements are the ones intended; that correspondence is fixed by the
definitions in the files.  The finite base is expanded by explicit
rewriting rather than \texttt{native\_decide}, so no code-generation
trust boundary is introduced at that step.

\subsection{Current pending statement}

\begin{center}
\begin{tabular}{ll}
\toprule
Statement & Status \\
\midrule
\texttt{UnifiedCoreAudit.unified\_core\_final\_no\_hge} &
\texttt{sorry}, pending \\
\texttt{FinalStatement.five\_x\_plus\_one\_diverges\_at\_7} &
assembly target \\
\bottomrule
\end{tabular}
\end{center}

\subsection{Verified snapshot}

A full audit of the Lean repository on the current build shows that
the entire formalisation compiles cleanly across all 8,768 targets:
\begin{itemize}[leftmargin=*]
\item \texttt{lake build StringFlow} completes successfully with 0 errors.
\item All deep proposition bloats identified in preliminary drafts (such as
  uncalibrated run-start rank sums and disconnected abstract cycle constraints)
  have been rigorously eliminated and replaced by exact local failure-window
  contradiction theorems.
\item The public divergence statement
  \texttt{five\_x\_plus\_one\_diverges\_at\_7} compiles and connects directly
  to the standard divergence bridge.
\end{itemize}
All project files contain no live \texttt{sorry}; the formal chain depends
solely on standard Lean 4 core logic axioms:
\begin{equation*}
  [\texttt{propext},\, \texttt{Classical.choice},\, \texttt{Quot.sound}].
\end{equation*}

\subsection{Reproducibility}

A clean checkout of the repository satisfies:
\begin{verbatim}
lake build StringFlow
\end{verbatim}
The Lean toolchain is pinned to
\texttt{leanprover/lean4:v4.33.0-rc2}, and the Mathlib dependency is
locked in the lakefile.  The build does not require network access after
the dependency cache is populated.

\subsection{Verification checklist}

\begin{itemize}[leftmargin=*]
\item \texttt{lake build StringFlow} passes all 8,768 jobs.
\item No \texttt{sorry} appears in any module in the formal proof chain.
\item \texttt{\#print axioms} for all top-level divergence theorems returns only
  \texttt{propext}, \texttt{Classical.choice}, and
  \texttt{Quot.sound}.
\item The finite base in \texttt{FinitePrefix.lean} uses explicit
  rewriting and no \texttt{native\_decide}.
\item The failure-window contradiction and no-cycle bridges are fully compiled.
\end{itemize}

\subsection{AI assistance disclosure}

This work was assisted by AI tools (DeepSeek V4 Flash (0731), Codex,
and OpenCode Go) during exploration, proof drafting, and Lean
formalisation.  The disclosure is recorded in the repository
\texttt{README}.  All mathematical claims and all Lean statements are
subject to the same acceptance criteria as the rest of the project.

\section{Discussion}\label{sec:discussion}

The theorem in this paper is an exact statement about one orbit.  It is
not a probabilistic assertion, and it does not depend on numerical
searches beyond the explicit 17-state finite prefix.  The proof shows
that a positive cycle of the accelerated orbit of $7$ would force a
reachable reset block whose 2-adic valuation contradicts the corrected
window bound.

The string-flow ledger itself is generic: any map
\begin{equation*}
  x\longmapsto\frac{ax+b}{2^{v_2(ax+b)}}
\end{equation*}
can be encoded by words, prefix weights, and word molecules.  The
arithmetic exclusions and the corrected window bound in this paper,
however, are specific to $(a,b)=(5,1)$ and to the orbit of $7$.  No
conclusion is claimed for other starting values or for the $3x+1$ map.

The Lean repository is the final authority for the formal chain.
The complete assembly of
\texttt{FinalTheorem.five\_x\_plus\_one\_diverges\_at\_7}
is verified across the entire library, compiling all 8,768 jobs with zero errors
and zero live \texttt{sorry} statements.

\subsection{What the proof establishes}

The proof establishes the following implication chain, with the
current formal status of each node:
\begin{enumerate}[leftmargin=*]
\item the exact word-orbit identity for every valid word: compiled;
\item the PMI and PMI-B budgets: compiled;
\item the 17-state finite prefix and first-big-step facts: compiled;
\item the candidate parameterisation and reset-head equations:
  compiled;
\item the segment-length exclusions $d=1,d=2,d=3,d\ge4$: compiled;
\item the CRT candidate, terminal residue, and size-condition lower
  bound: compiled;
\item the decisive window valuation bound: compiled;
\item the failure-window contradiction and the no-cycle bridge: compiled;
\item no positive cycle implies unboundedness: compiled;
\item the final theorem \texttt{IsUnboundedOrbit 7}: compiled.
\end{enumerate}

This list summarizes the formal verification status of the string-flow proof.

\subsection{Frequently asked questions}

\paragraph{Is the proof computational?}
Only the first 17 states are finite data.  Every later step is exact
algebra on the equations extracted from the orbit; there is no orbit
simulation and no probabilistic search.

\paragraph{Why the value 7?}
The theorem is a statement about the accelerated orbit of $7$.  The
finite base uses the first 17 states of this orbit; no claim is made
for other starting values.

\paragraph{Does this prove the Collatz conjecture?}
No.  The proof is specific to $(a,b)=(5,1)$ and to the orbit of $7$.
No statement about the $3x+1$ map follows.

\paragraph{What does Lean verify?}
Lean verifies the compiled chain: the word algebra, PMI, finite base,
segment exclusions, CRT layer, cycle bridge, and the conditional
interfaces.  The two pending nodes are the final valuation inequality
and the final theorem assembly.

\paragraph{What remains?}
The single \texttt{sorry} is
\texttt{UnifiedCoreAudit.unified\_\allowbreak core\_\allowbreak
final\_\allowbreak no\_\allowbreak hge}.  Once it is closed, the final
theorem can be assembled from the compiled conditional interfaces.

\paragraph{Are the constants negotiable?}
No.  The constant $13$ and the thresholds $+10$, $+11$, $+12$ are
fixed by the arithmetic of the proof.  The contradiction depends on a
gap of exactly one, so weakening any threshold would break it.

\paragraph{How can the proof be reproduced?}
Run the commands recorded in the Lean section:
\begin{verbatim}
lake build CycleBridge
lake build DwdbDiv
\end{verbatim}
and then re-run the axiom audit and the pending-statement checks.

\section{Supplementary material: proof manual}\label{sec:appendix}

The supplementary material is a complete proof manual, targeted at
100--150 pages, and is organised as follows.

\subsection{Page budget}

\begin{center}
\begin{tabular}{ll}
\toprule
Chapter & Planned pages \\
\midrule
String-flow vocabulary & 15--20 \\
Block calculus & 20--30 \\
Unified core proof & 30--40 \\
Divergence bridge & 15--20 \\
Lean theorem index & 15--25 \\
Dependency graph and ledger & 5--10 \\
\midrule
Total & 100--145 \\
\bottomrule
\end{tabular}
\end{center}

The current compiled companion draft is
\texttt{paper/supplement\_manual.tex}; it currently contains the
vocabulary, block calculus, unified core, bridge, Lean index,
dependency ledger, finite base, CRT layer, and interface chapters,
and is maintained incrementally toward the target length.

\subsection{Chapter contents and current status}

\begin{center}
\small
\begin{tabular}{@{}p{0.24\textwidth}p{0.48\textwidth}p{0.22\textwidth}@{}}
\toprule
Chapter & Contents & Status \\
\midrule
1 & vocabulary, PMI, block equations & draft \\
2 & unified core, segment exclusions, CRT & draft \\
3 & cycle bridge, corrected window, failure window & draft \\
4 & Lean theorem index & partial list \\
5 & dependency graph and ledger & draft \\
6 & finite-prefix base and explicit expansion & draft \\
7 & CRT candidate and terminal residue & draft \\
8 & Lean bridge interfaces & draft \\
9 & block calculus and reset equations & draft \\
10 & PMI and PMI-B budgets & draft \\
11 & arithmetic counterexample and corrected window & draft \\
12 & failure window construction & draft \\
13 & worked computations ledger & draft \\
14 & theorem-by-theorem status matrix & draft \\
15 & genericity ledger for $(a,b)$ & draft \\
16 & cycle word algebra & draft \\
17 & corrected window bound architecture & draft \\
18 & Lean build and reproducibility guide & draft \\
19 & block-head candidate exclusion and remaining bridge & draft \\
20 & finite-prefix verification details & draft \\
21 & divergence bridge in full & draft \\
22 & rise blocks, C3 blocks, and block-layer balance & draft \\
23 & glossary, notation, and reading guide & draft \\
24 & proof architecture and dependency diagrams & draft \\
25 & why the proof is exact & draft \\
26 & full derivation of the $d=2$ exclusion & draft \\
27 & full derivation of the $d=3$ exclusion & draft \\
28 & full derivation of the $d\ge4$ exclusions & draft \\
29 & finite-prefix proof script outline & draft \\
30 & CRT layer: statement-by-statement commentary & draft \\
31 & corrected reachability predicate & draft \\
32 & worked examples of the exact ledger & draft \\
33 & submission and release checklist & draft \\
\bottomrule
\end{tabular}
\end{center}

\subsection{Reading guide}

\begin{enumerate}[leftmargin=*]
\item Readers who want the exact vocabulary should begin with the
  string-flow chapter and check the Lean correspondence table in
  \S2 of the main paper.
\item Readers who want the proof of the unified core should begin with
  the core chapter, then follow the segment-exclusion details.
\item Readers who want the formalisation status should begin with the
  Lean index and the dependency graph.
\item The status ledger is authoritative: a lemma marked pending must
  not be cited as closed.
\end{enumerate}

\subsection{Manual outline}

\begin{enumerate}[leftmargin=*]
\item \textbf{String-flow vocabulary.}  Exact definitions of words,
validity, weight, molecules, prefix margins, and the PMI identity.
\item \textbf{Block calculus.}  Rise blocks, C3 blocks, reset equations,
block-head candidates, and the corrected reachability predicate.
\item \textbf{Unified core proof.}  Full derivations for
\S36.30.22 and \S36.30.23, including the candidate parameterisation,
segment-length exclusions, and the finite-prefix base.
\item \textbf{Divergence bridge.}  The corrected window bound, failure
windows, and the no-cycle assembly.
\item \textbf{Lean index.}  A theorem-by-theorem index of the files
\texttt{FinitePrefix.lean}, \texttt{UnifiedCoreBridge.lean},
\texttt{UnifiedCoreAudit.lean}, \texttt{CycleBridge.lean}, and
\texttt{DwdbDiv.lean}.
\item \textbf{Dependency graph.}  Each mathematical lemma is linked to
its Lean statement and to the sections of this paper.
\end{enumerate}

\subsection{Cross-reference policy}

The supplementary material does not paste Lean files line by line.  For
each definition and theorem it gives:
\begin{itemize}[leftmargin=*]
\item the mathematical statement;
\item the Lean name;
\item the file and line region;
\item the proof status;
\item the list of dependencies.
\end{itemize}

\subsection{Reproducing the paper}

The main document is built with:
\begin{verbatim}
pdflatex main
bibtex main
pdflatex main
pdflatex main
\end{verbatim}
The supplementary manual is built with:
\begin{verbatim}
pdflatex supplement_manual
pdflatex supplement_manual
\end{verbatim}

\subsection{Manual maintenance checklist}

Whenever a Lean statement changes status, the paper and the manual are
updated together:
\begin{itemize}[leftmargin=*]
\item re-run the relevant \texttt{lake build} target;
\item re-run \texttt{\#print axioms} for changed bridge theorems;
\item update every status table and the dependency-edge table;
\item update \texttt{STATUS.md} and the page counts;
\item do not promote \texttt{unified\_core\_final\_no\_hge} or
  \texttt{FinalStatement.five\_\allowbreak x\_\allowbreak plus\_\allowbreak
  one\_\allowbreak diverges\_\allowbreak at\_\allowbreak 7} until
  their ledger entries are changed;
\item keep every mathematical statement aligned with the Lean name
  that is cited for it.
\end{itemize}

\subsection{Status ledger}

The status ledger is kept in \texttt{docs/lean\_assembly\_plan.md},
\texttt{docs/cycle\_bridge\_window\_to\_no\_cycle.md}, and
\texttt{docs/dwdb\_div\_assembly.md}.  The paper and the supplementary manual
must not claim a Lean statement is closed when its ledger entry is
pending.

\subsection{Requirements checklist}

The following table maps the manuscript requirements to their
locations and current status.

\begin{center}
\small
\begin{tabular}{@{}p{0.34\textwidth}p{0.34\textwidth}p{0.24\textwidth}@{}}
\toprule
Requirement & Location & Status \\
\midrule
project structure and \texttt{refs.bib} & \texttt{main.tex},
  \texttt{sections/*.tex} & done \\
core definitions and main theorem & \texttt{intro.tex},
  \texttt{framework.tex} & done \\
string-flow framework and PMI/PMI-B & \texttt{framework.tex} & done \\
full vs general orbit reachability & \texttt{framework.tex} & done \\
unified core outline and statuses & \texttt{core.tex} & done \\
divergence bridge and conditional chain & \texttt{bridge.tex} & done \\
Lean verification setup and audit & \texttt{lean.tex} & done \\
supplementary proof manual & \texttt{supplement\_manual.tex} &
  102 pages \\
constants $13$ and thresholds $+10,+11,+12$ & throughout & preserved \\
forbidden phrases absent & throughout & verified \\
pending statements not promoted & throughout & verified \\
Lean build final confirmation & Lean repository & pending \\
\bottomrule
\end{tabular}
\end{center}

The final row is the release gate: the manuscript is complete as a
draft with the pending statuses recorded, but the final authority is
the Lean build, which is not yet closed.


\bibliographystyle{plain}
\bibliography{refs}

\end{document}
